% --- Document Class ---
\documentclass[12pt,a4paper]{report}

% --- Font Encoding (required for accented characters: ü, ş, etc.) ---
\usepackage[T1]{fontenc}

% --- Times Text Font (pdflatex compatible) ---
\usepackage{newtxtext,newtxmath}

% --- Page Geometry ---
\usepackage[a4paper, top=2.5cm, bottom=2.5cm, left=3cm, right=2.5cm, headheight=15pt]{geometry}

% --- float package: lets a table/figure be pinned with [H] so it never
%     gets bumped onto its own "float page". Isolated float pages get
%     vertically centered by LaTeX with large rubber space above/below,
%     which is what made Table 3.1 look like it filled a whole page. ---
\usepackage{float}

% --- Text Positioning for Cover Page ---
\usepackage[absolute]{textpos}
\usepackage{graphicx}
\setlength{\TPHorizModule}{1pt}
\setlength{\TPVertModule}{1pt}
\textblockorigin{0pt}{0pt}

% --- Line Spacing (Single line as per VGU requirement) ---
\usepackage{setspace}
\singlespacing

% --- Page Numbering & Headers/Footers ---
\usepackage{seqsplit}
\usepackage{fancyhdr}
\pagestyle{fancy}
\fancyhf{}
\fancyhead[L]{\nouppercase{\leftmark}}
\fancyfoot[C]{\thepage}
\renewcommand{\headrulewidth}{0.4pt}
\renewcommand{\chaptermark}[1]{\markboth{\chaptername\ \thechapter.\ #1}{}}
\fancypagestyle{plain}{%
  \fancyhf{}%
  \fancyfoot[C]{\thepage}%
  \renewcommand{\headrulewidth}{0pt}%
}
\fancypagestyle{preliminary}{%
  \fancyhf{}%
  \fancyfoot[C]{\thepage}%
  \renewcommand{\headrulewidth}{0pt}%
}
\fancypagestyle{empty}{%
  \fancyhf{}%
  \renewcommand{\headrulewidth}{0pt}%
}

% --- Table of Contents ---
\usepackage[nottoc,notlof,notlot]{tocbibind}

% --- Chapter Heading Style (compact "N.  Title" format, matching VGU reference) ---
\usepackage{titlesec}
\titleformat{\chapter}[hang]
  {\normalfont\Large\bfseries}{\thechapter.}{10pt}{\Large\bfseries}
\titlespacing*{\chapter}{0pt}{0pt}{20pt}
% Keep the same compact style for unnumbered chapters (Declaration, Abstract, Bibliography, etc.)
\titleformat{name=\chapter,numberless}[hang]
  {\normalfont\Large\bfseries}{}{0pt}{\Large\bfseries}
\titlespacing*{name=\chapter,numberless}{0pt}{0pt}{20pt}

% --- TikZ for org chart ---
\usepackage{tikz}
\usetikzlibrary{positioning, shapes.geometric, fit}
% Externalize disabled - not well supported on Overleaf
% \usepgfplotslibrary{external}
% \tikzexternalize
% \tikzset{external/system call={pdflatex -interaction=batchmode -halt-on-error -shell-escape #1}}

% --- Tables ---
\usepackage{booktabs}
\usepackage{array}
\usepackage[table]{xcolor}
\usepackage{float}
\usepackage{longtable}

% --- Code listings ---
\usepackage{listings}
\usepackage{caption}
\usepackage{amsmath}
\usepackage[backend=biber, style=apa, sorting=nyt]{biblatex}
\DeclareLanguageMapping{english}{english-apa}
\addbibresource{references.bib}

% --- Numbered APA-style bibliography (matches VGU reference thesis format) ---
\usepackage{enumitem}
\defbibenvironment{bibliography}
  {\list
     {[\arabic{enumiv}]}
     {\setlength{\leftmargin}{1.8em}%
      \setlength{\itemindent}{0pt}%
      \setlength{\labelwidth}{1.6em}%
      \setlength{\labelsep}{0.2em}%
      \setlength{\itemsep}{\baselineskip}%
      \setlength{\parsep}{\bibparsep}%
      \usecounter{enumiv}}}
  {\endlist}
  {\item}

% Colors matching a syntax-highlighted editor theme
\definecolor{codebg}{gray}{0.97}
\definecolor{codeframe}{gray}{0.55}
\definecolor{codekeyword}{RGB}{0,0,190}
\definecolor{codecomment}{RGB}{0,128,0}
\definecolor{codestring}{RGB}{163,21,21}
\definecolor{codenumber}{gray}{0.5}

\lstset{
  backgroundcolor=\color{codebg},
  basicstyle=\ttfamily\small,
  keywordstyle=\color{codekeyword}\bfseries,
  commentstyle=\color{codecomment}\itshape,
  stringstyle=\color{codestring},
  numbers=left,
  numberstyle=\tiny\color{codenumber},
  numbersep=10pt,
  showstringspaces=false,
  tabsize=4,
  frame=single,
  rulecolor=\color{codeframe},
  frameround=tftf,
  xleftmargin=2em,
  framexleftmargin=1.8em,
  breaklines=true,
  breakatwhitespace=false,
  captionpos=b,
}

% Numbers listings caption (italic)---explicit setup so listings 1.11b registers the counter.
\captionsetup[lstlisting]{labelfont={bf,small},textfont={it,small},labelsep=period}
\AtBeginDocument{\counterwithin{lstlisting}{chapter}}

% --- Hyperlinks ---
% MUST precede hyperref so the [hyphens] option propagates into hyperref's
% internal url package load, enabling line breaks after '-' in paths.
\PassOptionsToPackage{hyphens}{url}
\usepackage[colorlinks=true, linkcolor=black, citecolor=black, urlcolor=black]{hyperref}
\usepackage{url}

% --- Paragraph formatting (VGU standard: indentation, no parskip) ---
\usepackage{indentfirst}
\setlength{\parindent}{0.7cm}

% --- Semantic code formatting macros ---
% Usage: \code{tsc --noEmit}, \pkg{@hookform/resolvers}, \filename{src/database/schema/},
%        \apipath{'use client'} or \code{server-only}, \sqlkw{SELECT}
% \code uses \detokenize so '--' inside shell commands stays literal (no en-dash).
% All typewriter macros route through \nolinkurl (verbatim, with [hyphens] option
% from the url package loaded earlier) so long identifiers and paths can break
% at /, ., _, -, or alphanumeric boundaries instead of overflowing the column.
% \detokenize is preserved for \code to keep shell tokens literal.
\newcommand{\code}[1]{\texttt{\nolinkurl{\detokenize{#1}}}}
\newcommand{\pkg}[1]{\texttt{\nolinkurl{#1}}}
% \filename uses \nolinkurl (verbatim, no hyperlink) with \urlstyle{tt} to keep
% monospace appearance while enabling line breaks at /, ., _, and -.  Braces
% in path arguments are rendered verbatim; use <name> syntax for placeholders.
\DeclareRobustCommand{\filename}[1]{\bgroup\urlstyle{tt}\nolinkurl{#1}\egroup}
\newcommand{\apipath}[1]{\texttt{\nolinkurl{#1}}}
\newcommand{\sqlkw}[1]{\texttt{\nolinkurl{#1}}}

% ============================================
% DOCUMENT CONTENT
% ============================================
\begin{document}

% ============================================
% --- PRELIMINARY PAGES (Roman numerals) ---
% ============================================

% Cover Page (no page number)
\thispagestyle{empty}
\input{thesis_cover.tex}
\clearpage

\pagenumbering{roman}
\pagestyle{fancy}

\chapter*{Declaration}
\addcontentsline{toc}{chapter}{Declaration}

I, Do Hoang Trieu, hereby declare that this thesis, entitled ``Design and Pilot Protocol of a Next.js Boilerplate Framework for AI-Assisted SaaS Development,'' is my own original work. All sources of information have been properly acknowledged and referenced.

In preparing this thesis, I used AI tools to assist with code generation, prose drafting, and information research. I reviewed the output against the original sources and the SmartKit codebase, and take full responsibility for the content.

This thesis has not been submitted, in whole or in part, for the award of any other academic degree at this or any other institution. The author agrees that this thesis, upon acceptance, will be archived in the library of Vietnamese--German University and made available for reference use in accordance with the university's regulations.


\clearpage

\chapter*{Acknowledgements}
\addcontentsline{toc}{chapter}{Acknowledgements}

This thesis represents the capstone of my Bachelor's degree, a milestone that
would not have been achieved without the invaluable support of many individuals.
Therefore, I would like to express my deepest appreciation to those who have
accompanied me throughout this journey.

\textbf{My family} has been the cornerstone of this achievement. Their
unwavering encouragement, countless sacrifices, and steadfast belief in my
abilities carried me through every stage of this work. Without their quiet
strength and constant support, none of what follows would have been possible.

I am profoundly grateful to my supervisors for their dedicated mentorship.
My first supervisor, \textbf{Dr.~Le~Lam~Son}, provided exceptional guidance,
boundless patience, and steady encouragement throughout this research. His
rigorous academic oversight, incisive feedback, and continuous support were
invaluable in sharpening the technical depth of this thesis. Beyond this
project, I remain deeply grateful for his unwavering commitment, inspiring
leadership, and the profoundly positive influence he continues to exert on
our student community.

I am equally indebted to my second supervisor, \textbf{Dr.~Tran~Huu~Tam},
whose mentorship was instrumental in transforming a vague idea into a coherent
piece of engineering work. His dedication and encouragement helped shape this
project from its earliest stages into a structured thesis.

I would also like to thank the \textbf{Faculty of Engineering} and the
\textbf{Department of Computer Science} at the \textbf{Vietnamese--German
University} for laying a solid academic foundation over the past four years
and for fostering the environment in which this work could take shape.

Finally, my thanks go to my \textbf{university peers and friends}, whose
companionship through the shared struggles and memorable moments of our
undergraduate years made this journey not only bearable, but truly meaningful.

\clearpage

\chapter*{Abstract}
\addcontentsline{toc}{chapter}{Abstract}

AI coding assistants such as GitHub Copilot and Cursor AI are now in daily use, with substantial reported gains for boilerplate-heavy tasks: Peng et al.\ report a 55.8\% completion-time reduction on a controlled HTTP-server task \parencite{peng2023impact}, and Yetistiren et al.\ find boilerplate generation the strongest use case among practitioners \parencite{yetistiren2022assessing}. The benchmarks that dominate the literature, including HumanEval, MBPP+, APPS, BigCodeBench, and LiveCodeBench for algorithms and SWE-bench for bug fixing, measure problem-solving and patch generation, not whether AI-generated code respects framework conventions. Server/Client boundary violations in Next.js App Router, for instance, fail at runtime in a way no static type check catches.

This thesis proposes a three-dimensional evaluation framework for AI-generated Next.js boilerplate and instantiates it in SmartKit, an open-source Next.js 16.2.12 boilerplate that acts as both the framework's target codebase and a usable SaaS scaffold. The framework scores each generated feature on Structural Integrity, Type \& Contract Safety, and Coverage \& Completeness, each reported on a 0--5 scale. A pre-registered adaptive pilot ran 8 times on one feature: $n=5$ in Condition~A (SmartKit) after a mid-study extension, $n=3$ in Condition~B (convention-naive scaffold). Two of three directional predictions are confirmed outright; the third is asymmetric, confirmed in Condition~A and refuted in Condition~B, and is reported as such rather than collapsed into a single verdict.

The pilot is feasibility-stage evidence rather than an inferential test. Per-arm samples are too small for inferential statistics, and three protocol adaptations are disclosed: an adaptive extension applied to Condition~A only, a tool substitution for runs $6$--$8$, and two-clause reporting of the asymmetric prediction introduced after data collection. Larger-N replication, cross-tool validation, and developer-participant studies remain future work.

\clearpage

% Table of Contents
\tableofcontents

\clearpage

% List of Figures
\listoffigures
\thispagestyle{preliminary}

\clearpage

% List of Tables
\listoftables
\thispagestyle{preliminary}

\clearpage


% ============================================
% --- MAIN CONTENT (Arabic numerals) ---
% ============================================
\pagenumbering{arabic}
\pagestyle{fancy}

% ============================================
% CHAPTERS
% ============================================

% --- Chapter 1: Introduction and Motivation ---
\chapter{Introduction and Motivation}
\label{sec:introduction}

The chapter opens with the recent wave of AI coding assistants and the parallel shift in Next.js toward the App Router, then identifies the absence of an evaluation framework for framework-specific convention adherence. From that gap the chapter derives the research objectives, questions, and hypotheses that structure the rest of the thesis, and closes with the scope within which the SmartKit framework is designed and evaluated.

\section{Overview}

Controlled studies report substantial productivity gains from AI coding assistants. Peng et al. showed that developers using GitHub Copilot completed tasks 55.8\% faster than control groups \parencite{peng2023impact}, while Yetistiren et al. found AI tools particularly effective for boilerplate-generation tasks \parencite{yetistiren2022assessing}. Next.js has emerged as the dominant React meta-framework for production web applications, with its App Router architecture introduced in version 13 fundamentally changing application structure through React Server Components and explicit Server/Client boundaries \parencite{nextjsdocs}. The failure mode that motivates this thesis is concrete (Figure~\ref{fig:motivated-example}): when an AI coding assistant generates a Server Component that imports a Client Component and passes an event handler as a prop, Next.js produces the runtime error ``Event handlers cannot be passed to Client Component props''---a failure that \code{tsc --noEmit} cannot catch because it has no type-level representation. During SmartKit's development, recurring AI generation errors were observed firsthand in Next.js App Router code, including authentication security flaws, component boundary violations, and outdated API usage from older training data; these are development-time observations rather than measured prevalence figures. Without systematic evaluation, developers cannot assess whether AI-generated Next.js code adheres to framework conventions or will fail in production.

Existing benchmarks address algorithmic problem-solving and bug-fixing capabilities but not framework convention adherence. HumanEval evaluates Python code generation for standalone functions \parencite{chen2021evaluating}; SWE-bench assesses AI ability to resolve GitHub issues \parencite{jimenez2024swebench}; RigorBench (a recent, as-yet-unreviewed preprint, with a self-disclosed author conflict of interest) measures process-discipline metrics \parencite{rigorbench2026}. None of them capture the specific architectural requirements of Next.js App Router, such as Server/Client component boundaries, type-safe contracts, or file-system routing conventions. The gap means developers have no quantitative basis for comparing AI tools on framework-specific tasks.

\textbf{Intended users.} SmartKit targets solo developers and small teams (2--5 members) working on Next.js projects who use AI coding assistants (primarily Cursor AI and GitHub Copilot) for boilerplate generation. The intended users are intermediate developers who understand React and TypeScript fundamentals but seek to reduce boilerplate setup time and enforce consistent conventions in AI-assisted development. The framework supports incremental adoption: teams can apply the feature-based folder structure and Server/Client boundary conventions without committing to the full stack.

\section{The Problem: AI Convention Violations in Next.js}

This thesis addresses a concrete problem: AI coding assistants consistently generate Next.js App Router code that violates framework conventions, producing runtime errors and technical debt, yet no evaluation framework exists to measure this adherence. Designing such an evaluation framework faces three principal challenges. First, framework conventions are implicit and distributed across documentation, error messages, and community patterns, making them difficult to enumerate into testable metrics. Second, AI-generated code failures manifest at runtime rather than compile time. Server/Client boundary violations produce ``Event handlers cannot be passed to Client Component props'' errors only when executed, so they require dynamic analysis rather than static linting alone. Third, the evaluation dimensions must be designed to be orthogonal: structural integrity violations do not imply type errors, and coverage gaps operate independently of both, demanding a multi-dimensional framework rather than a single composite score.\footnote{The orthogonality of the three dimensions is asserted at design stage and awaits empirical validation; the pilot's per-arm sample is far below the $n \approx 30$ required to confirm that the dimensions vary independently in practice. See \S7.4 for the sample-size discussion that bounds this claim.}

The focus on Next.js App Router is deliberate. First, market dominance makes it the primary target: community surveys consistently place Next.js among the most widely adopted React meta-frameworks \parencite{nextjsdocs}; Stack Overflow's annual Developer Survey has tracked Next.js as the most-used React meta-framework since 2021 \parencite[Stack Overflow Developer Survey]{stackoverflow2025}. Second, training-data mismatch creates measurable failures: AI tools generate outdated patterns from older training data that conflict with current Next.js conventions, particularly around Server/Client boundaries and async API changes. Third, practical feasibility applies: the author has existing Next.js expertise, enabling direct testing of enforcement mechanisms. The pilot targets Cursor AI as the primary case study because its AGENTS.md feature implements documentation-based guidance, which prior research on grounded prompting suggests can improve AI output quality \parencite{barke2022grounded,zhang2023practices}.

\section{Research Objectives and Questions}

While Next.js dominates React meta-framework adoption and AI coding assistants such as GitHub Copilot and Cursor AI have demonstrated substantial productivity gains, there is a distinct gap at the intersection of these two areas: a systematic, operationalizable evaluation framework that measures how closely AI-generated Next.js boilerplate code adheres to framework-specific conventions. Existing benchmarks---HumanEval, SWE-bench, RigorBench---address algorithmic problem-solving and bug-fixing capability, but none capture the architectural requirements of Next.js App Router such as Server/Client component boundaries, type-safe contracts, or file-system routing conventions \parencite{chen2021evaluating,jimenez2024swebench,rigorbench2026}. This thesis targets precisely that gap.

The primary motivation of this thesis is to bridge this gap by designing a three-dimensional evaluation framework that turns an implicit set of framework conventions into measurable, comparable scores, and by instantiating that framework in SmartKit, a structured Next.js boilerplate that provides the necessary enforcement mechanisms and the evaluation target codebase.

The research pursues three objectives:
{\sloppy
\begin{itemize}
    \item \textbf{Design a three-dimensional evaluation framework.}
    Construct a framework that scores AI-generated Next.js boilerplate
    code along three dimensions: Structural Integrity, Type \& Contract Safety, and
    Coverage \& Completeness, so that convention adherence becomes a quantitative,
    comparable property of generated code.
    \item \textbf{Implement SmartKit as a structured architecture.}
    The thesis releases SmartKit, an open-source Next.js boilerplate
    that demonstrates feature-based modularity, explicit Server/Client
    boundaries, and type-safe contracts, providing both the enforcement
    mechanisms and the target codebase for the framework.
    \item \textbf{Conduct a pilot empirical case study.} The thesis applies the
    framework to a representative feature generated by Cursor AI under
    the SmartKit condition to demonstrate that the framework can
    produce illustrative, actionable scores on real generated code.
\end{itemize}
}

Consequently, this research seeks to answer the following research questions:
\begin{enumerate}
    \item How can an evaluation framework be consistently designed to
    measure AI coding assistant adherence to Next.js App Router boilerplate conventions along the Structural Integrity, Type \& Contract Safety, and
    Coverage \& Completeness dimensions?
    \item Can the resulting framework be operationalized into measurable
    metrics, and does it produce illustrative, actionable scores when
    applied to representative AI-generated Next.js features?
\end{enumerate}

These objectives and questions give rise to two classes of claims, both
documented as design-level pre-registrations to be exercised by the pilot
(Section~\ref{sec:pilot}). The first class states \textbf{framework
properties}---descriptive claims about whether the framework can measure
the dimensions it targets. The second class states \textbf{directional
predictions}---falsifiable claims about the sign of the score difference
between conditions, to be reported as confirmed, refuted, or inconclusive
per Section~\ref{sec:pilot}.

\paragraph{Framework properties.} The pilot is designed to demonstrate that
the framework can be operationalised on a real artifact:
\begin{itemize}
    \item \textbf{FP1 (Structural Integrity).} The framework can measure
    file placement, Server/Client boundary usage, and naming patterns in
    AI-generated Next.js code, producing a 0--5 score per run on a
    representative feature.
    \item \textbf{FP2 (Type \& Contract Safety).} The framework can measure
    TypeScript error counts and Zod-validation coverage in AI-generated
    Next.js code, producing a 0--5 score per run that combines a
    Type-Safety sub-score and a Contract-Safety sub-score via the
    weakest-link rule.
    \item \textbf{FP3 (Coverage \& Completeness).} The framework can
    enumerate AI-generated files against a feature-specific checklist and
    produce a 0--5 per-run score that reflects required-artifact presence.
\end{itemize}

\paragraph{Directional predictions.} Each directional prediction states
the expected sign of the difference between Condition~A (SmartKit at full
setup) and Condition~B (convention-naive scaffold) on a representative
feature---the \code{isEnabled(flagKey, userId?)} utility, which uses no
service that SmartKit integrates. Predictions are evaluated as
\emph{confirmed} (all runs in the predicted direction),
\emph{refuted} (no runs in the predicted direction), or
\emph{inconclusive} (intermediate), per the adaptive-N protocol
documented in Section~\ref{sec:pilot}.
\begin{itemize}
    \item \textbf{DP1 (Structural).} Condition~A is predicted to place the
    utility inside \filename{features/flags/lib/is-enabled.ts} in the
    majority of runs; Condition~B is predicted to place it at
    \filename{lib/feature-flags.ts} or at the repository root.
    \item \textbf{DP2 (Type \& Contract).} Condition~A is predicted to
    export a Zod-derived \code{FlagKey} enum or branded type in the
    majority of runs; Condition~B is predicted to use raw \code{string}
    keys.
    \item \textbf{DP3 (Coverage).} Condition~A is predicted to include a
    centralized \filename{src/tests/features/flags/is-enabled.test.ts} in the majority of runs;
    Condition~B is predicted to include no test file in the majority of
    runs.
\end{itemize}

Framework properties and directional predictions serve different
purposes. A framework property is a descriptive claim about the
framework itself; it cannot be refuted by pilot data, only supported or
unsupported. A directional prediction is a falsifiable claim about the
artifact and is the only class of claim to which the binary
\emph{confirmed/refuted/inconclusive} verdict applies. The pilot
protocol (Section~\ref{sec:pilot}) documents how each is reported. Full
statistical validation---including larger sample sizes, cross-tool
replication, and developer-participant studies---remains future
empirical work.

Designing an evaluation framework for AI-generated Next.js boilerplate encounters three key challenges that distinguish it from existing code generation benchmarks. First, framework conventions are implicit and distributed across documentation, error messages, and community patterns, requiring deliberate extraction into enumerable, testable rules rather than relying on predefined quality criteria. Second, convention violations in Next.js App Router manifest primarily at runtime. Server/Client boundary violations produce errors only when executed, so evaluation must go beyond static analysis to capture dynamic execution behaviour. Third, the evaluation dimensions are orthogonal: structural integrity violations do not imply type errors, and coverage gaps operate independently of both, demanding a multi-dimensional framework rather than a single composite score.\footnote{As in \S1.3 above, this orthogonality claim is asserted at design stage and awaits empirical validation; see \S7.4 for the sample-size discussion that bounds it and the planned Pearson correlation analysis that would distinguish genuine orthogonality from subtle coupling.}

The contributions of this thesis are as follows. First, it proposes a three-dimensional evaluation framework (Structural Integrity, Type \& Contract Safety, and Coverage \& Completeness) for measuring how closely AI coding assistants adhere to Next.js App Router conventions. The framework addresses a gap left by existing benchmarks, which score algorithmic correctness rather than framework-specific quality. Second, it releases SmartKit, an open-source Next.js boilerplate demonstrating feature-based modularity, explicit Server/Client boundaries, and type-safe contracts. SmartKit supplies both the conventions the framework measures and the target codebase it measures them on; all five Better Auth tables (user, session, account, verification), the billing feature (SePay integration with HMAC signature verification and idempotent webhook handling), and the email integration (Resend + React Email) are implemented at the current snapshot rather than deferred to future phases, supporting SmartKit's role as a functional SaaS boilerplate. Third, it releases the pilot protocol and the baseline comparison design as artefacts: a pre-registration document, a Condition~A vs.\ Condition~B comparison scaffold, a scoring rulebook, and a baseline-bootstrap script. The artefacts are concrete enough that another researcher can replicate the pilot without re-deriving the protocol from the thesis prose.

\begin{figure}[H]
  \centering
  \includegraphics[width=0.88\textwidth]{figs/fig-motivated-example.pdf}
  \caption{AI-generated Next.js App Router code violates Server/Client boundaries by passing event handlers across component types, producing runtime errors that static analysis cannot detect.}
  \label{fig:motivated-example}
\end{figure}

Figure~\ref{fig:evaluation-framework} summarizes the resulting evaluation-framework architecture. Each of the three dimensions is paired with an automated measurement: Structural Integrity is measured through parsed-AST rule violations, Type \& Contract Safety through TypeScript compiler output, and Coverage through a rule-coverage matrix checked against the convention catalog. The three measurements are applied to a shared corpus of AI-generated artifacts and produce a per-tool violation matrix, type-error rates, and coverage scores. Because the dimensions are deliberately independent, no single composite score can substitute for examining all three independently.

\begin{figure}[H]
  \centering
  \includegraphics[width=0.88\textwidth]{figs/fig-evaluation-framework.pdf}
  \caption{Three orthogonal evaluation dimensions (structural integrity, type \& contract safety, coverage \& completeness) applied to a target corpus of $5 \times 3 \times 10 = 150$ AI-generated artifacts; the executed pilot exercises $1 \times 3 \times 8 = 8$ runs.}
  \label{fig:evaluation-framework}
\end{figure}

This corpus scale represents the target scope for full empirical validation (Section~\ref{sec:future-work}); the thesis pilot (Section~\ref{sec:pilot}) applies the framework to a single representative feature within this scope.

\section{Scope and Limitations}

To keep the empirical rigor and feasibility of this research within the resource envelope of a single-author thesis, the scope is strictly defined by the following parameters and limitations:

\begin{enumerate}
    \item \textbf{Framework scope:} Next.js App Router only. No comparison with Nuxt, Remix, or other meta-frameworks.
    \item \textbf{AI assistants:} Cursor AI used for the pilot empirical case study. Broader AI assistant evaluation (GitHub Copilot, Claude Code) is future work.
    \item \textbf{Project type:} Single representative Next.js project demonstrating SaaS patterns.
    \item \textbf{Pilot study limitation:} The empirical case study is a worked-example executed against a total of $8$ runs under a pre-registered adaptive schedule, distributed as $n=5$ Cond~A and $n=3$ Cond~B on a single representative feature, the \code{isEnabled(flagKey, userId?)} utility, which uses no service that SmartKit integrates. The pilot illustrates framework applicability and produces per-condition medians plus directional predictions (DP1--DP3, Section~\ref{sec:pilot}); two of three are confirmed outright and the third is split between conditions. Cond~B remains at $n=3$, so the per-arm sample is still far below the $n \approx 30$ required for inferential validation (Section~\ref{sec:future-work}). Cross-tool replication, larger-$N$ studies, and developer-participant validation remain future empirical work.
\end{enumerate}

These boundaries are deliberate methodological choices. Fixing the framework, the AI assistant, and the project type removes confounding variance so that any measured difference between structured and unstructured conditions can be attributed to the SmartKit approach alone.


\chapter{Background}
\label{sec:background}

Understanding the intersection of AI coding assistants and Next.js development requires grounding in three domains: the capabilities and limitations of current AI coding tools, the architectural patterns of the Next.js framework, and the principles of SaaS boilerplate design.

\section{AI Coding Assistants}

GitHub Copilot, Anthropic's Claude, and Cursor AI all use transformer-based LLMs fine-tuned for code understanding and generation. GitHub Copilot provides inline code completions and entire function generation within supported editors (using OpenAI's GPT series and Anthropic's Claude series models, with specific model versions varying over time). Anthropic's Claude offers conversational code assistance through the Claude Code CLI and IDE integrations. Cursor combines Claude with a custom codebase indexing system that enables context-aware suggestions across entire projects.

These assistants share common architectural elements. They process code through transformer-based models trained on large corpora of publicly available source code, learning statistical patterns that inform code generation. When producing completions, they consider surrounding code context, natural language comments, and project-wide information to generate relevant suggestions. Their effectiveness varies with programming language, task complexity, and code structure \parencite{yetistiren2023evaluating}.

Research documents both strengths and limitations. Dakhel et al. found that AI assistance benefits expert developers who can filter suboptimal suggestions but may harm novice developers who lack the expertise to identify incorrect or inefficient code \parencite{dakhel2023github}. Liu et al. documented that ChatGPT performs better on older problems (pre-2021) than newer ones, with a 48.14\% advantage on the Accepted rate, though it struggles with complex algorithms and domain-specific code, frequently generating solutions with security vulnerabilities and inefficient implementations \parencite{liu2024chatgpt}.

The task taxonomy for AI suitability shows that boilerplate generation represents a strength area. Zhang et al. observed that developers primarily use Copilot for boilerplate generation, test writing, and code explanation rather than novel algorithmic development \parencite{zhang2023practices}. Pandey et al. confirmed this pattern, finding strong AI performance on template generation and documentation tasks but weak performance on complex logic and integration code \parencite{pandey2024transforming}. These findings focus the study on boilerplate interaction as an appropriate domain for AI enhancement.

\section{Next.js Framework Architecture}

Next.js is a React meta-framework that provides server-side rendering, static site generation, and API capabilities within a unified application model. Version 13 introduced the App Router paradigm, fundamentally changing how React applications structure pages, layouts, and data fetching. Developers who learned the Pages Router encounter several App Router concepts that the older model did not surface, including Server/Client boundaries and file-system conventions for special files.

React Server Components represent a core innovation in Next.js App Router. These components execute exclusively on the server, which lets them access databases and sensitive operations directly without exposing logic to the client. The distinction between Server and Client components requires explicit marking through the \apipath{'use client'} directive, creating a boundary that developers must carefully manage. Incorrect placement of components across this boundary produces runtime errors and performance degradation.

The App Router organizes applications through a file-system-based routing mechanism. Each folder within the \filename{app/} directory corresponds to a route segment, with \filename{page.tsx} files defining the content for each route. Special files such as \filename{layout.tsx}, \filename{loading.tsx}, and \filename{error.tsx} provide shared UI and error handling across route hierarchies. This convention-driven structure enables rapid scaffolding but requires adherence to specific naming and placement patterns.

TypeScript integration in Next.js emphasizes type safety throughout the application. The framework provides typed interfaces for page properties, API request and response objects, and data fetching results. Achieving full type coverage requires consistent use of TypeScript patterns including generics, utility types, and strict mode configuration. AI coding assistants may generate code that compiles but violates framework conventions, producing runtime errors or subtle bugs.

\section{SaaS Boilerplate Patterns}

Software-as-a-service applications share common requirements that boilerplate frameworks address: user authentication, database access patterns, multi-tenancy support, and deployment configurations. The economics of SaaS development favor approaches that reduce per-project initialization effort while maintaining code quality and consistency.

Krüger et al. established that platform reuse reduces maintenance effort but involves upfront investment in framework adoption and learning curve \parencite{kruger2020empirical}. The trade-off explains the appeal of well-designed boilerplate templates that amortize initialization cost across multiple projects. Aleem et al. identified customization, scalability, multi-tenancy, security, integration, and fault tolerance as key factors in SaaS architecture decisions \parencite{aleem2019maturity}.

Modern scaffolding approaches build on model-driven architecture principles. Inayatullah and Azam showed that automated code generation reduces boilerplate development effort for cross-platform applications \parencite{inayatullah2019model}. Liu and Sra recently showed that LLMs have potential for automatically generating scaffolded user interfaces, reducing complexity for users in professional creative software \parencite{liu2025auggen}.

Meta-framework evaluation identifies architectural trade-offs across options. Prior work on framework adoption in software engineering has identified learnability, extensibility, and ease of migration as primary drivers of adoption, while framework lock-in and high configurability are commonly cited as major adoption risks \parencite{storey2016frameworkadoption}. Next.js App Router emphasizes React Server Components as a first-class primitive, while maintaining full compatibility with Client Components for interactive UI. The hybrid model creates specific conventions that boilerplate frameworks should encode.

\chapter{Related Work}
\label{sec:relatedwork}

This section surveys existing research on code generation benchmarks, empirical studies of AI-assisted development, and systematic literature reviews on AI for software engineering. The analysis identifies the specific gap this thesis addresses: the absence of framework-specific evaluation for AI coding assistants operating on Next.js boilerplate conventions.

\section{Code Generation Benchmarks}

The evaluation of AI coding assistants relies heavily on standardized benchmarks that measure code generation capability. Chen et al. introduced HumanEval as a foundational benchmark containing 164 hand-written programming problems with function signatures, docstrings, and unit tests \parencite{chen2021evaluating}. The benchmark evaluates models using Pass@1, Pass@10, and Pass@100 metrics. OpenAI's Codex achieved 28.7\% Pass@1 on this benchmark, establishing a baseline for subsequent model development.

HumanEval focuses exclusively on single-function Python code generation with algorithmic and reasoning requirements. The design prioritizes general programming capability over domain-specific or framework-aware generation. Yu et al. extended the benchmark concept with HumanEval Pro, adding self-invoking code generation tasks that require models to generate functions calling previously defined functions \parencite{yu2025humanevalpro}. Their work documented up to 30\% performance degradation on self-invoking tasks compared to standalone tasks, indicating that current benchmarks underestimate real-world code generation complexity.

SWE-bench represents a substantial advance in benchmark realism by evaluating AI models on real GitHub issues from popular Python repositories \parencite{jimenez2024swebench}. Each task requires understanding the issue description, locating relevant code, implementing a fix, and passing unit tests. Jimenez et al. found that state-of-the-art models resolved only a small fraction of issues---Claude 2 solved approximately 1.96\% and GPT-4 only 1.7\%---and the gap between benchmark performance and practical software engineering capability is substantial.

Language-specific benchmarks identify further evaluation limitations. Petrukha et al. developed SwiftEval for Swift programming, finding that general-purpose benchmarks overestimated Swift code generation capability \parencite{petrukha2025swifteval}. Their language-specific benchmark documented 40--60\% performance gaps not captured by cross-language evaluations.

No existing benchmark evaluates framework convention adherence, boilerplate structure quality, TypeScript or Next.js-specific patterns, or multi-file scaffolding tasks. The benchmarks measure whether AI assistants can solve algorithmic problems or fix bugs in existing codebases, not whether they can generate framework-compliant boilerplate from scratch.

Several extensions and alternatives warrant mention. Liu et al. developed HumanEval+ and MBPP+ as rigorous extensions of the original benchmarks, finding that pass rates drop often by 30\% or more when evaluated against additional high-quality tests per problem \parencite{liu2024correct}. Zhuo et al. introduced BigCodeBench, which evaluates 1,140 fine-grained tasks combining diverse function calls with complex natural-language instructions \parencite{zhuo2024bigcodebench}. Hendrycks et al. proposed APPS, a 10,000-problem benchmark spanning introductory, interview, and competition-level coding challenges \parencite{hendrycks2021apps}. Jain et al. addressed the critical contamination problem through LiveCodeBench, which sources problems from recent competitive programming contests published after model training cutoffs \parencite{jain2024livecodebench}. Li et al. developed CodeReviewer for code review automation, an adjacent evaluation direction that assesses code quality dimensions relevant to AI-assisted development workflows \parencite{li2022codereviewer}. Each of these benchmarks measures distinct aspects of code generation or code quality capability, yet none evaluate the framework-specific convention adherence that the SmartKit framework targets.

\section{Empirical Studies on Developer Productivity}

Several empirical studies examine how AI assistance affects real developer workflows. Peng et al. conducted the seminal controlled experiment on GitHub Copilot with 95 recruited developers implementing an HTTP server in JavaScript \parencite{peng2023impact}. Their methodology used a treatment-control design with matched tasks, finding the 55.8\% completion time improvement cited earlier. Heterogeneous effects showed that less experienced developers benefited most from AI assistance.

The enterprise-level .NET study by Tariq et al. extended this methodology to enterprise software development scenarios \parencite{tariq2026enterprise}. With 60 developers and 960 task observations across eight enterprise-oriented tasks, they measured task completion time, implementation errors, cyclomatic complexity, maintainability index, code smells, and NASA-TLX cognitive load scores. AI-assisted development reduced task completion time, errors, code smells, and cognitive workload while improving maintainability. The trade-off was that AI generated redundant abstractions and increased code verbosity, with effectiveness depending on developer expertise and continued oversight.

Liu et al. assessed ChatGPT code generation quality across multiple languages and task types \parencite{liu2024chatgpt}. ChatGPT performs better on pre-2021 problems with a 48.14\% advantage on Accepted rate, but its ability to directly fix erroneous code through multi-round dialogue is relatively weak. Generated code often contained security vulnerabilities and inefficient implementations requiring human verification, though ChatGPT could resolve over 89\% of identified vulnerabilities through multi-round fixing. Yetiştiren et al. documented that Copilot generates functionally correct code for many problems but frequently produces solutions with higher cyclomatic complexity and lower readability than human-written solutions \parencite{yetistiren2022assessing}.

AlOmar et al. examined developer-ChatGPT refactoring conversations, finding that developers frequently needed multiple iterations to achieve satisfactory refactoring \parencite{alomar2024refactoring}. AI often missed design-level improvements, and common anti-patterns emerged in developer-AI refactoring interactions. AI effectiveness therefore depends on the structure of the interaction and on the developer's ability to guide refinement, not only on the tool itself.

Hou et al. conducted a systematic literature review of 395 research articles on LLMs for software engineering from 2017--2024 \parencite{hou2024llmse}. Code generation (35\%) and code repair (25\%) are the most studied application areas. Most studies use benchmark evaluation (60\%) rather than human evaluation, and reported effectiveness varies widely by task type. The review identified research gaps in security, requirements engineering, and empirical studies. Framework-specific boilerplate evaluation would begin to address these gaps.

Russo et al. analyzed GenAI adoption complexity, identifying five key dimensions: technical, organizational, individual, process-related, and ethical \parencite{russo2024navigating}. Technical capability was necessary but insufficient; human factors including trust, skill, and workflow integration were critical success factors. Evaluating framework-specific boilerplate tools therefore requires attention to how well they integrate with developer workflows and support skill development.

\section{Evaluation Methodology}

The refinement of AI-generated code has received attention as a complementary evaluation dimension. Liu et al. analyzed 4,066 ChatGPT-generated programs across 2,033 programming tasks in Java and Python \parencite{liu2024refining}. Of these programs, 2,756 were deemed correct, 1,082 produced wrong outputs, and 177 contained compilation or runtime errors; 1,930 additional programs suffered from maintainability issues identified through static analysis. Their automated refinement pipeline using static-analysis feedback improved code quality by more than 20\%. Evaluation frameworks should therefore measure both initial generation quality and the effort required for refinement.

Ouyang et al. studied the non-determinism of ChatGPT in code generation \parencite{ouyang2025nondeterminism}. Across 829 problems in three benchmarks, the ratio of coding tasks with zero equal test output across different requests ranged from 47.56\% (HumanEval) to 75.76\% (CodeContests). Setting the temperature to 0 does not guarantee determinism, despite being a common practice. Reliable assessment therefore requires multiple runs with confidence intervals rather than single-point measurements.

Beller et al. proposed the DAT framework for measuring software development productivity at Meta \parencite{beller2026dat}. The approach combines diff authoring time, survey-based metrics, and AI tool-specific productivity measures. The DAT framework models a path for evaluating AI-boilerplate interaction that balances quantitative performance metrics with qualitative developer experience assessment.

Nguyen-Duc et al. outlined a research agenda for GenAI in software engineering, emphasizing code generation quality, testing and verification, developer experience, education and training, ethics and security, and process and project management \parencite{nguyenduc2025genai}. The agenda explicitly calls for standardized evaluation and longitudinal studies, requirements that framework-specific evaluation frameworks can begin to satisfy.

Sauvola et al. projected four primary future scenarios for generative AI in software development, ranging from legacy integration to clean-slate and networked development, with software production shifting toward higher-level design and oversight roles as routine coding tasks become increasingly automated \parencite{sauvola2024future}. Research on optimizing AI interaction patterns, including framework-specific boilerplate conventions that reduce the cognitive overhead of AI-assisted development, follows from this projection.

\section{Research Gap Summary}

The reviewed literature establishes that AI coding assistants improve developer productivity on repetitive tasks, that boilerplate generation represents an AI strength area, and that existing benchmarks measure algorithmic capability rather than framework convention adherence. To the author's knowledge---given the structured-but-non-exhaustive literature search described in the next section---no prior work has specifically evaluated AI tool effectiveness on Next.js App Router conventions, React Server Components patterns, TypeScript-first boilerplate structures, or full-stack framework-specific best practices. This gap is independently supported by Hou et al.\ (2024)'s systematic review of 395 LLM-SE articles, which does not list framework-specific convention adherence as a separately tracked evaluation category \parencite{hou2024llmse}; the review's documented research gaps (security, requirements engineering, empirical studies) likewise do not overlap with framework-specific boilerplate evaluation. Together with RigorBench's complementary observation that existing benchmarks measure outcomes rather than process discipline \parencite{rigorbench2026}, this triangulates the gap claim against a peer-reviewed systematic review rather than relying on a single, COI-disclosed preprint. This thesis proposes a Next.js boilerplate framework engineered for AI interaction and an evaluation methodology that measures framework-specific quality dimensions, reporting both as design-stage contributions.

\section{Literature Search Methodology}

The literature reviewed in this chapter was identified through a structured search process. Primary databases searched include ACM Digital Library, IEEE Xplore, arXiv, and Google Scholar. Search terms combined AI coding assistant terminology (``LLM code generation,'' ``Copilot evaluation,'' ``ChatGPT code quality'') with framework-specific terminology (``React Server Components,'' ``Next.js,'' ``TypeScript boilerplate'') and benchmark terminology (``code generation benchmark,'' ``framework-specific evaluation''). Inclusion criteria prioritized peer-reviewed publications from 2020 onward, with seminal pre-2020 works (e.g., HumanEval, Codex) included for foundational context. The 35 references in the bibliography are supplemented by the systematic literature review of Hou et al. \parencite{hou2024llmse} which provides independent coverage of 395 articles in the LLM-for-SE space. The author acknowledges that a fully PRISMA-compliant systematic review with explicit screening protocol is beyond single-author-thesis scope; the search described here is structured but not exhaustive, and additional framework-specific evaluation work may exist outside the search scope.

\section{Alternative Approaches to AI-Codebase Alignment}

The structural approach taken by SmartKit is one of several strategies that researchers and practitioners have explored to improve AI-codebase alignment. This section briefly reviews these alternatives and justifies why the structural approach was selected for this thesis.

\textbf{Prompting strategies} represent the most lightweight alternative. Enhanced system prompts, few-shot examples, and retrieval-augmented generation (RAG) can improve AI output quality without modifying project structure. Prompting approaches suffer from context window limitations, require repeated re-instruction across sessions, and cannot enforce consistency across a codebase without external governance. RigorBench (2026) found that existing benchmarks evaluate outcomes but not the engineering process through which AI agents arrive at solutions \parencite{rigorbench2026}. Prompting alone cannot address this structural process-discipline gap.

\textbf{Enhanced documentation} such as thorough READMEs, inline comments, and architecture decision records (ADRs) can guide AI behavior. Documentation requires human maintenance, is easily ignored by AI systems generating new code, and does not provide machine-readable enforcement mechanisms.

\textbf{Linter and formatter configurations} (ESLint, Prettier) can catch some convention violations but operate post-generation rather than preventing them. These tools validate style, not semantics; they cannot enforce feature boundaries or Server/Client component conventions that require architectural awareness.

The \textbf{structural approach} taken by SmartKit addresses these limitations by encoding conventions at the project structure level. Conventions become discoverable by AI systems through file location and naming patterns, are enforced through project structure rather than developer vigilance, and provide reusable templates that scale across features. The structural approach works alongside prompting and documentation, not as a replacement.

\paragraph{Caveat on the alternative-approaches dismissals.} The dismissals of prompting, enhanced documentation, and linter-based enforcement above are \emph{design-stage reasoning} based on prior literature, not empirical comparisons. The pilot's Condition~A advantage is the \emph{combined} effect of multiple overlapping signals---the \filename{AGENTS.md} instructions, the visible example features (\filename{features/auth/}, \filename{features/billing/}), and the existing linter/Zod configuration---and the executed pilot does not isolate ``structural'' conventions from the prompting, documentation, and tooling layers that operate alongside them. A controlled comparison of structural vs.\ prompting vs.\ documentation vs.\ linter-only conditions is left as future work (Section~\ref{sec:future-work}); the present thesis reports only the bundled effect of the SmartKit convention envelope.

Table~\ref{tbl:starter-comparison} positions SmartKit qualitatively against five established Next.js starter templates on seven dimensions relevant to AI-assisted boilerplate generation. The comparison is descriptive rather than empirical: cells summarise publicly documented features of each template, not measured outcomes, and no performance claim should be drawn from them.

\begin{table}[H]
\centering
\footnotesize
\renewcommand{\arraystretch}{1.3}
\setlength{\tabcolsep}{2.5pt}
\caption{Qualitative positioning of SmartKit against established Next.js starter templates. Cells describe documented features, not measured performance; cell-level claims were transcribed from each template's own documentation.\label{tbl:starter-comparison}}
\begin{tabular}{@{}p{2.5cm}*{6}{>{\centering\arraybackslash}p{1.95cm}}@{}}
\toprule
 & \textbf{SmartKit} & \textbf{create-t3-app} & \textbf{Next.js official} & \textbf{Vercel templates} & \textbf{shadcn/ui} & \textbf{bulletproof-react} \\
\midrule
Feature-based layout & Yes & Partial & No & No & No & Yes \\
Server/Client convention & \texttt{'use client'} & tRPC types & Manual & Manual & Manual & Manual \\
ORM included & Drizzle & Prisma & None & None & Drizzle & Prisma \\
Auth library & \pkg{better-}\pkg{auth} & NextAuth.js & None & NextAuth.js & Auth.js & Context API \\
TypeScript-first & Yes & Yes & Optional & Optional & Yes & Yes \\
Testing scaffold & Vitest & Vitest & None & None & Vitest & Vitest \\
AGENTS.md & Yes & No & No & No & No & No \\
\bottomrule
\end{tabular}
\renewcommand{\arraystretch}{1}
\end{table}

\footnote{\texttt{create-t3-app}: \url{https://create.t3.gg/}; Next.js official: \url{https://nextjs.org/docs/app/getting-started}; Vercel templates: \url{https://vercel.com/templates}; \texttt{shadcn/ui}: \url{https://ui.shadcn.com/}; \texttt{bulletproof-react}: \url{https://github.com/alan2207/bulletproof-react}. SmartKit claims trace to the source repository accompanying this thesis.}

The qualitative positioning in Table~\ref{tbl:starter-comparison} suggests a structural approach to convention encoding in AI-assisted development. While most frameworks provide partial feature-based layout or manual Server/Client conventions, SmartKit includes \filename{AGENTS.md} as a first-class project artifact---enabling persistent, non-negotiable convention encoding directly readable by AI coding assistants. This positioning is asserted as design-stage reasoning and is not an empirical claim; cell-level differences across rows are descriptive notes rather than verified performance gaps.

\chapter{Proposed Solution}
\label{sec:solution}

This section presents the design of SmartKit, a Next.js boilerplate framework engineered to optimize interaction with AI coding assistants. The design synthesizes empirical findings from the literature with domain-specific analysis of Next.js conventions to produce a codebase structure that improves AI comprehension and output quality.

\section{Design Rationale}

Zhang et al. found that structured, well-documented codebases improve AI coding assistant performance \parencite{zhang2023practices}. Barke et al. showed that grounded prompting, providing relevant code context, improves AI output quality \parencite{barke2022grounded}. The SmartKit design applies these findings through concrete architectural decisions: encode conventions explicitly so that AI systems can discover them, group related code so that context stays narrow, and define type-safe contracts so that data shapes remain explicit.

These principles translate into three concrete design choices. Explicit convention encoding makes framework-specific patterns visible and discoverable by AI systems. Feature-based modularity groups related code to enable context-aware suggestions. Type-safe contracts between components give AI systems reliable information about data structures and function signatures.

\section{Feature-Based Modular Architecture}

Many Next.js projects organize code by technical layer (components/, lib/, hooks/) or by page structure. SmartKit instead adopts feature-based modularity, where each feature contains all code necessary for its functionality. An authentication feature, for instance, includes database schemas, API routes, React components, form validation schemas, and email templates within a single directory.

This organization improves AI comprehension by narrowing the search space for relevant context when AI generates code for a specific feature. The scope of changes becomes explicit, which reduces the risk of AI-generated code violating feature encapsulation. The organization also creates natural boundaries for testing and deployment that align with developer mental models.

The feature directory structure follows a consistent pattern across all features:

\begin{lstlisting}[
  language={},
  caption={Feature directory layout convention used across SmartKit.},
  label={lst:feature-layout},
  basicstyle=\ttfamily\small,
  showstringspaces=false,
]
src/features/{feature-name}/
+-- components/         # React components
+-- lib/               # Feature-specific utilities
+-- types/             # TypeScript type definitions
+-- emails/            # Email templates
+-- index.ts          # Public API barrel
\end{lstlisting}

This structure allows AI systems to predict where to place new code within a feature, reducing the cognitive overhead of project organization.

\section{Server/Client Boundary Conventions}

React Server Components introduce an explicit boundary between server-only and client-accessible code. SmartKit uses Next.js's native Server/Client boundary mechanism through the \apipath{'use client'} directive for Client Components and the \pkg{server-only} import for server-only utilities. The pattern extends Next.js's convention-based routing to component architecture.

Server/Client boundaries affect data flow design. Server Components can access databases and environment variables directly, while Client Components must receive data through props or server actions. AI systems can invoke type-safe server actions from Client Components, which keeps
data flow consistent as new components are generated.

Listing~\ref{lst:agents-md} shows a representative excerpt from the SmartKit \filename{AGENTS.md} file, which encodes the conventions that Cursor AI reads when operating inside the project. The file defines project-level rules that take precedence over instructions found in code comments or tool outputs, and includes specific conventions for file placement, Server/Client boundaries, and schema sharing.

\begin{lstlisting}[
  language={},
  caption={Representative excerpt from SmartKit's AGENTS.md file.},
  label={lst:agents-md},
  basicstyle=\ttfamily\small,
  showstringspaces=false,
]
# AGENTS.md

Project-level rules that AI agents must follow when working
inside the SmartKit repository. These rules are non-negotiable.

## Thesis Formatting
- Citations & References -- Provide references for all content.
- Source Code Listings -- Present as a listing (not a figure).
- Figures & Tables Numbering -- All must be explicitly referred to.

## Convention Rules
- Feature directories: place all feature code under src/features/{name}/
- Client Components: use 'use client' directive
- Server-only utilities: use server-only import
- Zod schemas: shared between client validation and server-side parsing
- Server actions: src/features/{name}/actions/
\end{lstlisting}

\section{Database Schema Patterns}

SmartKit uses Drizzle ORM with PostgreSQL to provide type-safe database access throughout the application. The schema design follows conventions that AI systems can learn and reproduce. Each table schema includes columns with explicit types, relations defined with foreign keys, and indexes specified for query optimization.

The Drizzle schema pattern enables AI-generated code to access the database through type-safe query builders. AI systems can examine existing schema patterns and generate consistent queries that maintain type safety across the application. Ad-hoc database access that bypasses type checking or violates query conventions is discouraged.

Schema files for Better Auth's shared tables (\code{user}, \code{session}, \code{account}, \code{verification}) are all implemented in \filename{src/database/schema/}, an infrastructure layer outside the FSD hierarchy, alongside the billing-specific tables (\code{billing}, \code{payment\_transactions}) used by \filename{features/billing/}. The five Better Auth tables cover the \code{user} identity, the \code{session} and \code{account} records used for sign-in flows, and the \code{verification} table for email-verification challenges; relations between them are declared in \filename{src/database/schema/relations.ts}. All Drizzle schema definitions are kept in this single directory so that \pkg{drizzle-kit} can scan them reliably and all features can import from a common location. The placement keeps database patterns centralized while remaining directly importable from any feature module.

\section{Type-Safe API Contracts}

API routes in SmartKit use Zod schemas for request and response validation. These schemas serve dual purposes: runtime validation of incoming requests and TypeScript type inference for handler code. AI systems can examine Zod schemas to understand API contracts, generating code that respects type constraints without requiring explicit type annotations.

The shared schema pattern enables consistent validation across frontend form components and backend API handlers. The same Zod schemas validate both client-submitted forms and server-handled requests. The match between frontend and backend validation reduces the gap between AI-generated client-side and server-side code, producing more coherent implementations.

\chapter{Implementation}
\label{sec:implementation}

This section describes the technical implementation of the SmartKit framework using Next.js 16.2.12 and associated technologies. The implementation shows how the design principles translate into working code.

\section{Technology Stack}

Each component of the stack was chosen because it reinforces one of the three evaluation dimensions rather than for its own sake. Next.js 16.2.12 supplies the App Router architecture that Structural Integrity measures directly (Server/Client boundaries, file-system routing), with TypeScript~5.7.3 in strict mode providing the compile-time signal that Type~\&~Contract Safety reads from \code{tsc --noEmit}. Drizzle~ORM~0.45.2 (PostgreSQL) and Zod~4.4.3 jointly enforce the framework's type-safe-contract principle (Section~\ref{sec:solution}): Zod schemas are colocated in \filename{features/<name>/types/}, consumed by Server Actions for runtime parsing, and re-used by client-side forms via \pkg{z.infer}, so a single schema is the source of truth on both sides of the Server/Client boundary that Coverage~\&~Completeness checks for. Server Actions in \filename{features/<name>/actions/} handle mutations and call \code{revalidatePath} after writes, letting Server Components fetch data directly through Drizzle and pass typed props to Client Components without an intermediate client-side data layer.

The remaining tooling is chosen for consistency rather than for a direct role in the rubric: Tailwind~CSS~4.3.3 (CSS-first configuration) for styling, Biome~2.5.6 as a single linter/formatter in place of separate ESLint and Prettier configurations, Resend~6.18.1 with React~Email~6.9.1 for transactional email templates, and SePay as the payment gateway integrated through \filename{features/billing/} (HMAC signature verification + idempotent webhook handling; described in detail in \S\ref{sec:billing-sepay}). \pkg{@better-auth/drizzle-adapter} integrates authentication with the database layer. None of these four are exercised by the pilot feature (Section~\ref{sec:integrations-and-pilot}); they matter for SmartKit's second role as a usable SaaS starting point rather than for the evaluation itself.

\section{Project Structure}

The SmartKit project structure implements the feature-based architecture described in Section~\ref{sec:solution}:

\begin{lstlisting}[
  language={},
  caption={Top-level directory structure of the SmartKit instantiation.},
  label={lst:project-structure},
  basicstyle=\ttfamily\small,
  showstringspaces=false,
]
smartkit/
+-- src/
|   +-- app/              # Next.js App Router pages
|   +-- features/         # Feature-based modules
|   |   +-- auth/         # Authentication feature
|   |       +-- components/   # React components
|   |       +-- emails/       # Email templates
|   |       +-- lib/          # Auth-specific helpers
|   |       +-- types/        # Auth type definitions
|   |       +-- index.ts      # Public API barrel
|   +-- database/         # Drizzle schema, migrations, client
|   |   +-- schema/       # user, session, account, verification (Better Auth tables) + billing + payment_transactions
|   |   +-- migrations/   # Drizzle migration files
|   |   +-- db.ts         # Drizzle client
|   +-- lib/              # Shared utilities
|   +-- config/           # Site and environment config
+-- docs/                 # Documentation
\end{lstlisting}

The \filename{app/} directory contains Next.js App Router pages and layouts. Route-specific code lives here, while feature code lives in \filename{src/features/}. The separation keeps page composition concerns distinct from feature implementation.

The \filename{database/} directory contains Drizzle configuration and schema definitions. It hosts the database schema---all five Better Auth tables (\filename{user.ts}, \filename{session.ts}, \filename{account.ts}, \filename{verification.ts}, plus the relations declared in \filename{relations.ts}) and the billing-specific tables (\filename{billing.ts} with \code{payment\_transactions})---in \filename{database/schema/}. Keeping all schemas in one location enables \pkg{drizzle-kit} to scan consistently and allows any feature module to import shared types without cross-feature directory dependencies.

The \filename{lib/} directory contains shared utilities that span multiple features, such as date formatting helpers, class name utilities, and common type definitions.

The \filename{features/flags/} directory at the repository top-level (a deliberate placement outside of \filename{src/}, mirroring the AGENTS.md convention \emph{up to the \filename{src/} prefix}) is a representative minimal feature added by the pilot's baseline-bootstrap script to provide a non-trivial integration target for the framework's Coverage dimension. The top-level placement was chosen so that the AI scanning the repo root encounters \filename{features/flags/} immediately when searching for a feature-flag implementation, rather than descending into \filename{src/} first; this is consistent with the AGENTS.md rule that \emph{new} feature code should be placed under a feature-named directory but does not enforce the \filename{src/} prefix for the pilot's representative feature. It follows the same internal conventions as \filename{src/features/auth/} (described in detail below): \filename{lib/} for the \code{isEnabled} implementation, \filename{types/} for the Zod-derived \code{FlagKey} schema, and \filename{index.ts} as the public API barrel, with the corresponding unit test at \filename{src/tests/features/flags/is-enabled.test.ts} (which \emph{is} inside \filename{src/}, in line with the centralized-test convention documented in the AGENTS.md file). Its presence is what makes the pilot executable against the framework without having to first create a feature by hand.

\section{Key Implementation Patterns}

The Server/Client boundary convention introduced in Section~\ref{sec:solution} is enforced consistently at implementation level: every generated file is checked for the \apipath{'use client'} directive or \pkg{server-only} import appropriate to its role, which is what \texttt{score-flag.ts} inspects for the Structural Integrity dimension.

Database queries use Drizzle's type-safe query builder with explicit schema references. Generated code can use TypeScript's type inference to maintain consistency with existing queries. The query patterns include pagination, filtering, and relation loading that AI systems can adapt for new features.

API routes use Next.js Route Handlers with Zod schema validation. Request handlers receive parsed and validated data through a consistent interface, while response formatting uses a standard structure that includes status codes, data payloads, and error messages.

\section{Billing and SePay Integration}
\label{sec:billing-sepay}

The \filename{features/billing/} feature implements a payment and order-management flow backed by SePay, a Vietnamese payment gateway. This feature is the second-tier pattern source visible to the AI in Condition~A alongside \filename{features/auth/}, so its directory layout and conventions are worth describing explicitly.

The feature follows the same feature-based layout used elsewhere: \filename{actions/} for Server Actions, \filename{components/} for React components, \filename{hooks/} for client-side hooks, \filename{lib/} for business logic and external integrations, \filename{types.ts} for shared type definitions, and \filename{index.ts} as the public API barrel. The top-level \filename{features/billing/index.ts} re-exports the public surface (entitlements helpers, error types, content templates) and keeps internal helpers private.

\textbf{SePay webhook handling.} The webhook entry point is \filename{process-sepay-payload.ts}, which validates the inbound payload, looks up the corresponding order, and routes the result to either order completion or a recorded failure. Signature verification lives in \filename{signature.ts}, which performs HMAC-SHA256 verification of the \code{X-Sepay-Signature} header against the shared secret stored in environment variables. The verification helper returns a discriminated result (\code{ok} | \code{mismatch} | \code{malformed}) so callers do not need to inspect raw header values, and the ordering helper is unit-tested against canonical SePay fixtures.

\textbf{Idempotency.} Webhook delivery can be retried by SePay under transient failures, so the handler must deduplicate. Idempotency is provided through the \filename{payment-transactions.ts} module, which records each processed transaction id keyed by the gateway-supplied reference. A repeated reference is short-circuited as a no-op rather than re-running the order-completion side effects. The same module exposes a query helper used by reconciliation jobs (\filename{cancel-expired-orders.ts}, \filename{process-order-completion.ts}) to detect stuck orders.

\textbf{Database tables.} The billing feature is backed by the \filename{billing} table family in \filename{src/database/schema/billing.ts}, which defines the order, transaction, and entitlement tables; \filename{payment-transactions} is the dedup key for the idempotency check described above. The \filename{account}, \filename{session}, \filename{user}, and \filename{verification} tables (Better Auth) are joined on \code{userId} to attach payment state to authenticated identities; these are listed under \filename{src/database/schema/} and are all implemented at the current snapshot rather than deferred to future phases.

This feature is not exercised by the pilot (\S\ref{sec:integrations-and-pilot}); it is mentioned here because its directory layout, signature verification pattern, and idempotency discipline are what the AI pattern-matches against in Condition~A, and because the boilerplate's value as a working SaaS starting point depends on these production-grade integrations being present rather than aspirational.

\chapter{Evaluation}
\label{sec:evaluation}

This section presents the proposed evaluation framework and pilot feasibility study for assessing AI coding assistant performance on Next.js boilerplate generation tasks. The framework design is accompanied by a pilot evaluation protocol that exercises the framework on a representative feature.

\section{Evaluation Framework}
\label{sec:evaluation-framework}

The evaluation uses a multi-dimensional quality framework that captures different aspects of boilerplate quality. The framework extends existing benchmark methodology to address framework-specific concerns that existing evaluations overlook.

Recent systematic research has identified a fundamental gap in how AI coding benchmarks measure capability. RigorBench (2026), surveying nine major AI coding benchmarks spanning function-level, issue-level, and project-level evaluation, found that none incorporate any measurement of engineering process discipline\parencite{rigorbench2026}. RigorBench proposes a process-oriented evaluation framework targeting the agent's workflow rather than project-specific convention adherence. The SmartKit framework is complementary: it measures convention adherence in generated boilerplate code, an output-quality dimension that RigorBench does not directly evaluate. The three-dimensional evaluation framework proposed in this thesis follows from this complementarity.

The three dimensions capture distinct, non-overlapping aspects of boilerplate quality that affect real-world usability. Rather than measuring only output correctness, the framework assesses whether generated code can be integrated into production codebases without convention violations or type safety issues. Each dimension maps to an automated measurement approach and corresponds to a specific failure mode observed in AI-generated Next.js code.

\textbf{Unit of analysis.} Scoring is performed per generated file: each file receives one score per dimension, and feature-level scores are computed as the median across constituent files. The median is used rather than the mean because individual files within a feature may exhibit extreme outlier scores due to edge-case generation failures (e.g., a single malformed file), and the median is robust to such outliers.

\textbf{Structural Integrity} assesses whether generated code follows Next.js conventions for file location, naming, Server/Client component boundaries, and App Router patterns. Automated checking uses the \texttt{score-flag.ts} script (file-placement checks, \texttt{server-only} import detection, \texttt{'use client'} directive checks, barrel-file verification) specific to Next.js structural conventions. The dimension captures violations such as incorrect file placement in the \filename{features/} directory, missing \apipath{'use client'} directives on client components, and non-standard naming patterns that break routing.

\textbf{Type \& Contract Safety} captures two related but distinct quality signals, treated as sub-scores that combine into a single Type \& Contract Safety score on the 0--5 rubric. \emph{Type Safety sub-score} focuses on compile-time checks: \code{tsc --noEmit} error counts and \pkg{@typescript-eslint/}\pkg{no-explicit-any} occurrences signal whether generated code introduces type regressions into the codebase. \emph{Contract Safety sub-score} focuses on runtime validation: the proportion of API route handlers that include Zod schemas for request/response validation indicates whether the generated feature is production-ready out of the box. Sub-score combination rule: the Type \& Contract Safety dimension score equals the lower of (Type Safety sub-score, Contract Safety sub-score) on the 0--5 scale, reflecting that a single weak sub-signal is not offset by a strong other. The combination rule is asserted at design stage; future empirical work should validate that ``weakest-link'' aggregation tracks developer experience better than weighted average or geometric mean alternatives.

\textbf{Coverage \& Completeness} evaluates whether generated boilerplate includes all necessary artifacts for a functional feature. Table~\ref{tbl:coverage-checklist} shows the reference checklist for the \code{isEnabled} feature flag utility used in the pilot. Completeness is measured by enumerating generated files against the checklist and computing a coverage ratio (files present divided by required files). Missing artifacts indicate incomplete boilerplate that requires additional developer effort to make functional.

\subsubsection{Relationship between SmartKit Integrations and the Pilot Feature}
\label{sec:integrations-and-pilot}

SmartKit instantiates the framework conventions across multiple feature areas, including billing (SePay integration), authentication (Better~Auth), and email (Resend). These integrations exist because the artifact is a working SaaS boilerplate, not because the evaluation framework requires them. The pilot deliberately selects a feature---the feature-flag check utility---that uses \emph{none} of the integrated services, in order to isolate the effect of convention encoding from the effect of having integrated services available. If the pilot used a feature that integrates SePay or Resend, any observed difference between Condition~A and Condition~B would conflate convention encoding with integration availability (\emph{confounding~B} in the discussion notes). The integrations remain first-class in the artifact; the pilot simply does not exercise them. This separation is what allows the framework to remain generalisable beyond SmartKit: a developer using a different stack can encode the same conventions without SePay, Resend, or Better~Auth, and the framework still measures the same three dimensions on any generated feature.

A second, related point concerns the role of the integrations during the AI session. In Condition~A, the AI sees the full SmartKit file tree, including \filename{features/billing/} (SePay HMAC, idempotency), \filename{features/auth/} (Better~Auth schemas and server actions), and the Resend email templates under \filename{features/auth/emails/}. These files are not used by the pilot feature, but they encode the SmartKit convention envelope: feature-based folder structure, colocated server actions, Zod schemas in \filename{types/}, \pkg{server-only} boundaries, and \code{revalidatePath} after mutations. The AI in Condition~A pattern-matches this envelope when generating the pilot feature, even though the pilot feature itself does not call any of the integrated services. In Condition~B, none of these pattern sources are available, so the AI must invent both the utility and its surrounding convention. The score difference between conditions therefore measures the convention envelope, not the integration surface. The convention envelope is what the framework is designed to capture; the integration surface is what makes SmartKit a usable SaaS boilerplate. Both are required for the thesis, but they play distinct roles. The integration surface contributes to the second thesis contribution---the SmartKit artifact as a working SaaS boilerplate---rather than to the pilot; the pilot is the instrument that exercises the framework on AI-generated code, while the integrations are the evidence that the framework conventions are usable in production code.

The standardised prompt's first resolution-order rule---``Per-user DB override, only if \code{userId} is passed''---is implemented
in the pilot as a stubbed path: the generated \seqsplit{readUserOverride} helper returns \code{undefined} without querying any database layer, so the pilot feature genuinely touches neither Drizzle nor any external persistence. The stub keeps the resolution order observable in tests while leaving full DB-backed override as future work; this is consistent with the feature's design choice to avoid integrating any SmartKit service, including Drizzle, into the pilot scope. The stub implementation is itself what caps the pilot's Coverage rubric score at level~3 (``Adequate''), since rubric A.5.3 level~4 requires the per-user override path to be implemented (not stubbed); the corrected Coverage cells in Table~\ref{tbl:pilot-actual} and the correction footnote immediately following that table document this relationship.

\begin{table}[ht]
\centering
\footnotesize
\renewcommand{\arraystretch}{1.25}
\caption{Reference coverage checklist for the \code{isEnabled} feature flag utility.}
\label{tbl:coverage-checklist}
\begin{tabular}{@{}p{0.7cm}p{3.3cm}>{\raggedright\arraybackslash}p{9cm}@{}}
\toprule
\# & Artifact & Expected path \\
\midrule
1 & Core utility function & \filename{features/flags/lib/is-enabled.ts} \\
2 & Zod schema and FlagKey enum & \filename{features/flags/types/flags.ts} \\
3 & Unit test file & \filename{src/tests/features/flags/is-enabled.test.ts} \\
4 & Public API barrel & \filename{features/flags/index.ts} \\
5 & Cache implementation & Embedded in \filename{is-enabled.ts} (5-minute in-memory cache) \\
\bottomrule
\end{tabular}

\vspace{2pt}
\footnotesize\textit{Note:} The \code{isEnabled(flagKey, userId?)} utility resolves feature flags with three override layers (per-user DB override, per-flag environment variable, default value) and caches results for five minutes. This feature was chosen for the pilot because it uses no service that SmartKit integrates (no Resend, Better Auth, SePay, or Drizzle dependency beyond the schema), eliminating integration-pattern confounding.
\renewcommand{\arraystretch}{1}
\end{table}

\section{Design Rationale and Trade-offs}
\label{sec:design-rationale}

The evaluation framework design involves several deliberate choices that balance scope with feasibility.

\textbf{Dimension selection.} The three dimensions were designed to address a gap that RigorBench (2026) identifies in existing benchmarks: they measure algorithmic correctness and engineering process discipline but not convention adherence in generated boilerplate \parencite{rigorbench2026}. Each SmartKit dimension targets a distinct failure mode observed in AI-generated Next.js code: structural violations (file placement, Server/Client boundaries), type safety gaps, and missing boilerplate artifacts.

\textbf{Automated over manual evaluation.} The framework prioritizes automated metrics to enable reproducibility and scale. Expert review provides richer qualitative assessment, but automated checking ensures consistent evaluation across the large number of code samples required for meaningful comparison. Where qualitative judgment is necessary (e.g., structure quality assessment), LLM-as-judge with majority voting is a design-stage option that was not exercised in the executed pilot; the pilot uses single-primary-rater manual scoring (deferred IRR; see \S7.2).

\textbf{6-level rubric across dimensions.} Each of the three~dimensions is scored on a 0--5 rubric (six integer levels). The per-dimension cell descriptions are detailed in \S A.5.1--A.5.3; the rough mapping is: $0$ = missing artefact, $1$ = severe violation, $2$ = partial (in the right neighbourhood but wrong sub-detail), $3$ = acceptable, $4$ = good (level 3 + a specific bonus criterion, e.g.\ \code{server-only} for Structural), $5$ = excellent (level 4 + a second bonus criterion, e.g.\ barrel re-export). The exact level descriptions are dimension-specific and asserted at the design stage; empirical validation of the cutoff mapping to developer outcomes is future work (\S7.4). The aggregate per-run score is the sum of the three dimensions, range $0$--$15$.

\textbf{Single-framework focus.} The evaluation focuses exclusively on Next.js to maintain depth over breadth. Future work should extend the framework to Nuxt and Remix to assess generalizability of the approach.

\section{Evaluation Scenarios}
\label{sec:evaluation-scenarios}

The evaluation scenario measures AI-generated boilerplate quality directly: AI systems receive task descriptions, and the generated code is measured against the framework dimensions.

Scenario A provides AI systems with feature descriptions and requests complete feature scaffolding. The evaluation measures the completeness and quality of generated components, schemas, API routes, and tests. Future iterations of Scenario A would evaluate multiple AI systems (GitHub Copilot, Claude, Cursor AI) to enable comparative analysis; the current pilot focuses on Cursor AI as the primary case study, consistent with the scope defined in Section~\ref{sec:introduction}.

Future work should evaluate these scenarios with developer participants following the methodology of Peng et al. \parencite{peng2023impact} and Tariq et al. \parencite{tariq2026enterprise}.

All evaluations use multiple runs to address non-determinism concerns raised by Ouyang et al. \parencite{ouyang2025nondeterminism}. The Cursor AI tool does not allow user-settable temperature in Auto mode (the tool snapshot used in the pilot); multi-run evaluation therefore substitutes for temperature variation as the determinism-handling mechanism. Results report per-condition medians and ranges, plus Wilson 95\% confidence intervals on the binary pass rate ($\geq 3/5$).

\section{Control Comparison}
\label{sec:control-comparison}

Table~\ref{tbl:condition-ab} summarises the two conditions used in the
pilot evaluation. Condition~A is SmartKit at full setup; Condition~B is
a convention-naive Next.js scaffold produced by
\code{create-next-app} with the same dependency surface as SmartKit but
no FSD layout, no AGENTS.md, and no \filename{.cursor/rules/} folder.
The two conditions therefore differ only in convention encoding, which
is the variable the framework is designed to measure. The exact commit
and CLI invocation used to build each condition are recorded in
Appendix~\ref{app:pre-registration} for replication.

\begin{table}[H]
\centering
\small
\renewcommand{\arraystretch}{1.2}
\caption{Condition~A (SmartKit at full setup) vs.\ Condition~B
(convention-naive scaffold) used in the pilot evaluation.}
\label{tbl:condition-ab}
\begin{tabular}{@{}p{2.6cm}p{5.6cm}p{5.6cm}@{}}
\toprule
\textbf{Dimension} & \textbf{Condition~A (treatment)} & \textbf{Condition~B (control)} \\
\midrule
Stack surface
  & SmartKit (Next.js 16, Drizzle, Biome, Zod, Resend, Better~Auth)
  & create-next-app scaffold + full SmartKit dependency surface (Next.js 16, TypeScript, Tailwind CSS, Biome, Drizzle ORM, Zod 4.4.3, Resend, Better Auth); no \filename{src/features/} layout, no AGENTS.md, no \filename{.cursor/rules/} \\[4pt]
\midrule
Convention envelope
  & FSD \filename{src/features/<name>/}, AGENTS.md, \pkg{server-only}, colocated Zod in \filename{types/}, Server Actions in \filename{actions/}
  & No AGENTS.md, no \filename{src/features/}, flat \filename{lib/} at root \\[4pt]
\midrule
Folder layout
  & Feature-scoped: \filename{components/}, \filename{hooks/}, \filename{actions/}, \filename{lib/}, \filename{types/}, \filename{emails/} per feature
  & \filename{app/} (pages), \filename{lib/} (flat), \filename{components/} (flat) \\[4pt]
\midrule
AI tool context
  & Sees existing \filename{features/auth/}, \filename{features/billing/} patterns; conventions live in repo files, not in separate docs
  & Sees only \code{create-next-app} default plus env placeholders; no SmartKit conventions to read \\[4pt]
\midrule
Pilot feature
  & Generate \code{isEnabled(flagKey, userId?)} utility + Zod schema + centralized test at \filename{src/tests/features/flags/}
  & Generate same \code{isEnabled(flagKey, userId?)} utility; no SmartKit conventions to read \\[4pt]
\bottomrule
\end{tabular}
\renewcommand{\arraystretch}{1}
\end{table}

\subsection{Design Rationale}

A central threat to the validity of the pilot results is confounding: the structured SmartKit condition and the unstructured condition differ \emph{primarily} in convention encoding, but they also differ in second-order ways that the executed pilot does not fully control. The most important residual confound is that the SmartKit repository exposes several example features (\filename{features/auth/}, \filename{features/billing/}, \filename{features/auth/emails/}) that the AI can use as pattern-matching anchors beyond the explicit AGENTS.md instructions. To isolate SmartKit's contribution from this confound, the pilot compared it against an unstructured Next.js baseline sharing the same dependency surface as SmartKit (Next.js 16.2.12, TypeScript 5.7.3, Tailwind CSS 4.3.3, Drizzle ORM 0.45.2 with PostgreSQL, Zod 4.4.3, Biome 2.5.6, Resend 6.18.1, and \pkg{better-auth} 1.6.x, plus the \filename{src/database/schema/} directory inherited from SmartKit) but with no feature-based \filename{src/features/} layout and no AGENTS.md. The exact dependency surface and CLI invocation used to build the baseline are recorded in Appendix~\ref{app:pre-registration} (Conditions~\ref{app:cond-a}--\ref{app:cond-b}) for replication; the \code{smartkit-baseline/package.json} shipped alongside this thesis is the canonical record. The pilot feature (\code{isEnabled(flagKey, userId?)}) was chosen for the same reason: it uses no service that SmartKit integrates (no Resend, Better~Auth, SePay, or Drizzle dependency beyond the schema), so the AI cannot copy a pattern from the example features. This design follows established experimental practice for fair comparison \parencite{peng2023impact}. This section refers to the baseline as the ``convention-naive scaffold'' to avoid confusion with the SmartKit repository directory itself; the build script that produces it is preserved alongside the thesis artefacts (Appendix~\ref{app:pre-registration}).

The control artifact serves as the control condition, providing a baseline that isolates the effect of SmartKit's convention encoding from the effect of the underlying technology stack.

\subsection{Experiment Protocol}

The executed pilot applied the standardized prompt (Listing~\ref{lst:control-prompt}) verbatim to both codebases.

\begin{lstlisting}[
  language={},
  caption={Standardized prompt for the control comparison experiment, applied verbatim to both conditions.},
  label={lst:control-prompt},
  basicstyle=\ttfamily\footnotesize,
  showstringspaces=false,
  breaklines=true,
]
Add a feature flag check utility for this Next.js codebase. The
utility should expose a single function `isEnabled(flagKey, userId?)`
that returns a boolean.

Resolution order (highest priority first):
  1. Per-user DB override, only if `userId` is passed.
  2. Per-flag environment variable `FLAG_<KEY>=true|false`.
  3. Default value (passed as the second argument).

Constraints:
  - Cache the result in-memory for 5 minutes, keyed by flag and
    (if present) userId. Cache must be process-local.
  - The function should be a server-side utility; do not export it
    to the browser.
  - Define a `FlagKey` type that the function accepts. Use a Zod
    schema to derive it.
  - Include at least 3 unit tests covering: env override, default
    value, and cache hit.
  - Use TypeScript strict mode.

You may add files anywhere in the repository that you think is
appropriate. The only requirement is that the public API is reachable
as `isEnabled` and the type `FlagKey` is exported.
\end{lstlisting}

The prompt intentionally avoids referencing SmartKit-specific conventions (\emph{do not} say ``for SmartKit'', ``follow the feature-based layout'', ``use \filename{.server} naming'', or ``refer to AGENTS.md''). The same prompt applies to both SmartKit and the convention-naive scaffold, so any difference in output quality is attributable to the convention encoding rather than the prompt content. The full pre-registration record, including the frozen prompt hash, is reproduced in Appendix~\ref{app:pre-registration}.

The pilot followed a pre-registered adaptive schedule, detailed in Appendix~\ref{app:pre-registration}: the first three runs were executed in Condition~B, and the inter-run range on Type~\&~Contract Safety ($0$--$3$) exceeded the pre-registered continuation threshold, triggering Runs~4--5 in Condition~A. Condition~A was then extended from $n=2$ to $n=5$ (Runs~6--8) so the pre-registered binary verdict rule ($N=5$ per condition) could issue a determinate outcome; the extension and its justification are logged in the amendment history (Appendix~\ref{app:pre-registration}, Table~\ref{tab:amendment-log}). The final tally is eight runs: $n=5$ Condition~A and $n=3$ Condition~B. The same per-run pipeline was applied to all eight runs, with per-condition aggregates reported in Table~\ref{tbl:pilot-actual}:

\begin{enumerate}
    \item Initialize a fresh copy of the codebase with the pre-registered tooling verified against a run-day checklist.
    \item Invoke Cursor AI with the standardized prompt; start the time-to-correct clock.
    \item Collect generated files, Biome violations, \code{tsc --noEmit} output, and artifact enumeration; snapshot the run in \filename{runs/cond-<a,b>/run-<NN>/}.
    \item Score each run on the three-dimension rubric; record per-run JSON plus manual cells; commit before the next run.
\end{enumerate}

The control comparison produced a two-column table (SmartKit vs.\ convention-naive scaffold) showing dimensional scores per run, enabling direct attribution of any observed difference to the structural conventions.

\subsection{Expected Threats and Mitigations}

The control comparison addresses the primary threat to internal validity: that observed scores on SmartKit reflect technology stack quality rather than convention adherence. By using the same dependency surface in both conditions, the experiment isolates the convention variable. In the executed pilot, the Cond~B Type \& Contract dimension showed a wider inter-run range ($3$ points across $n=3$) than Cond~A; this residual within-condition variance---not convention encoding---is the most plausible alternative explanation for the size of the score gap, and is one reason the DP verdicts in Section~\ref{sec:pilot} are reported as descriptive rather than inferential. Cross-AI-tool replication, where version drift can be measured independently across tools, is identified as future work (Section~\ref{sec:future-work}).

\section{Multi-Run Evaluation}
\label{sec:multi-run}

\subsection{Non-Determinism Threat}

Ouyang et al. \parencite{ouyang2025nondeterminism} found that 47.56\% to 75.76\% of coding tasks produce non-deterministic outputs across requests on leading benchmarks, with temperature settings providing no reliable guarantee of determinism. Single-run evaluations therefore risk capturing idiosyncratic outputs rather than representative behavior. This threat is especially acute for framework-specific evaluation: an AI may generate convention-compliant code in one run and violate boundaries in the next, and a single-run score would not reflect this variance.

The evaluation addresses this threat through a multi-run protocol that records per-dimension scores per run and reports per-condition medians and ranges, plus Wilson 95\% confidence intervals on the binary pass rate ($\geq 3/5$) across multiple independent generation attempts.

\subsection{Experimental Protocol}

The multi-run protocol follows these steps:

\begin{enumerate}
    \item \textbf{Define temperature setting.} Record the temperature parameter (Cursor AI default: approximately 0.7).
    \item \textbf{Execute 5 independent runs.} For each run, clean the workspace of previously generated files before generating anew. Apply the same standardized prompt (Listing~\ref{lst:control-prompt}) to ensure consistency across runs.
    \item \textbf{Collect artifacts per run.} For each run, record: generated files snapshot, Biome violation counts, \code{tsc --noEmit} error counts, and artifact coverage ratio.
    \item \textbf{Score each run.} Apply the three-dimension rubric to each run independently, producing a per-run score on the 0--5 scale for each dimension.
    \item \textbf{Compute descriptive statistics.} For each dimension across the $n=5$ runs per condition:
    \begin{itemize}
        \item \textbf{Median score:} the middle value when scores are sorted, robust to outliers.
        \item \textbf{Range:} $[\text{min}, \text{max}]$ across the 5 runs, showing score variability.
        \item \textbf{Wilson score interval:} for each run treated as a binary outcome (pass = score $\geq 3/5$; fail = score $< 3/5$), compute the Wilson 95\% confidence interval on the proportion of passing runs. The Wilson interval is given by Equation~\ref{eq:wilson}:
        \begin{equation}
        \frac{k + \frac{z^2}{2}}{n + z^2} \pm \frac{z}{n + z^2} \sqrt{\frac{k(n-k)}{n} + \frac{z^2}{4}}
        \label{eq:wilson}
        \end{equation}
        where $k$ is the number of passing runs, $n=5$, and $z = 1.96$ for 95\% confidence. The Wilson interval is preferred over the normal approximation (Wald interval) for small $n$ because it never produces bounds outside $[0, 1]$ and has better coverage properties when the true proportion is near 0 or 1.
    \end{itemize}
\end{enumerate}

\subsection{Statistical Reporting}

Results are reported as median [min, max] with Wilson 95\% CI for the pass rate. For example: ``Structural Integrity scored median 3/5 [range: 2--4] across 5 runs; 4/5 runs passed the $\geq 3/5$ threshold, Wilson 95\% CI: [0.30, 0.99].'' A minimum detectable effect size should be pre-specified (e.g., 0.5 points on the 0--5 rubric) to guide sample-size planning for future extensions. If the Wilson CI lower bound exceeds the threshold of interest (e.g., 0.60 for ``most runs pass''), the framework can make a confident directional claim; if the CI spans 0.50, the claim is inconclusive pending additional runs.

\subsection{Implications for Framework Reliability}

The multi-run evaluation serves two purposes. First, it provides robustness evidence: if the 95\% CI for a dimension is narrow and consistently above a quality threshold, the framework produces reliable signals despite AI non-determinism. Second, it shows which dimensions are most stable across runs. Structural Integrity may show high variance if the AI places files inconsistently; Type \& Contract Safety may be more stable if Zod schema generation is driven by reliable TypeScript patterns. These patterns inform where framework enforcement mechanisms are most and least effective.

The pilot presented in Section~\ref{sec:pilot} ($n=5$ Cond~A, $n=3$ Cond~B after the pre-registered extension) was a feasibility-stage design selected for single-author-thesis scope. The multi-run protocol documented here remains the template for larger-$N$ extension, as it directly addresses the non-determinism threat documented in the published literature.

\section{Evaluation Metrics}

Quantitative metrics derive from automated analysis of generated code. \textit{Structural Integrity} uses Biome custom rules to check Next.js-specific file placement, naming, and Server/Client boundary conventions. \textit{Type \& Contract Safety} uses \code{tsc --noEmit} output to identify type violations and custom Biome rules for \emph{any}-type detection and Zod coverage. \textit{Coverage \& Completeness} uses file enumeration against a reference checklist per feature type.

Qualitative metrics derive from expert review of generated code. Reviewers assess whether generated code follows Next.js best practices, whether naming conventions are consistent, and whether the overall implementation aligns with framework expectations. The expert review layer follows an inter-rater reliability (IRR) protocol: two independent raters score the same generated artifacts, with a target Cohen's $\kappa$ $\geq$ 0.7. The automated scoring layer (Biome, \code{tsc --noEmit}) is deterministic and requires no IRR. The IRR protocol is documented here as a design-stage design; execution of the IRR pass is identified as future work pending availability of a second rater.

Productivity metrics such as diff authoring time and developer-reported productivity \parencite{beller2026dat} are not collected in the current pilot and remain future work; the executed pilot reports only per-run time-to-correct (ttc, in minutes) as a coarse timing signal. Task completion time, cyclomatic complexity, maintainability index, and code smell density are documented as design-stage metrics for a future developer-participant study.

\section{Pilot Evaluation Protocol}
\label{sec:pilot}

This section reports the pilot evaluation of the \code{isEnabled(flagKey, userId?)} utility under the pre-registered adaptive protocol, comparing Condition~A (SmartKit with full convention encoding) against Condition~B (convention-naive scaffold). The final tally is $8$ runs: $n=5$ Condition~A and $n=3$ Condition~B (Section~\ref{sec:control-comparison}). The pilot is reported as a worked example, not a hypothesis test: Two of the three directional predictions are confirmed outright (DP1, DP2); DP3 is confirmed on its Condition A clause and refuted on its Condition B clause under the pre-registered binary rule ($5/5$ Condition~A runs in the predicted direction), but the per-arm sample remains far below the $n \approx 30$ needed for inferential validation, so cross-tool replication, larger-$N$ validation, and developer-participant studies remain future work.

\textbf{Task.} The pilot applied the framework to a representative task: generating the \seqsplit{isEnabled(flagKey, userId?)} utility using Cursor AI. This utility resolves a feature flag with three override layers (per-user DB override, per-flag environment variable, default value) and caches the result in-memory for five minutes. As established in Section~\ref{sec:integrations-and-pilot}, the feature was chosen because it exercises none of SmartKit's integrated services, eliminating integration-pattern confounding. The task instruction mirrored a realistic developer request without over-specifying implementation details, so as to evaluate the AI assistant's ability to discover and apply existing conventions rather than merely follow explicit instructions.

\textbf{Scoring.} \textit{Structural Integrity} was scored by running the \texttt{score-flag.ts} script against the generated files, counting violations per category (file placement, naming, Server/Client boundary). \textit{Type \& Contract Safety} was scored by running \code{tsc --noEmit} on the generated code and counting type errors and \emph{any}-type occurrences. \textit{Coverage \& Completeness} was scored by enumerating generated files against a reference checklist (utility function, Zod schema, centralized test at \filename{src/tests/features/flags/}, public API barrel) and computing a coverage ratio. Each dimension was reported on the 0--5 rubric described in Sections~\ref{sec:evaluation-framework} and~\ref{sec:evaluation-scenarios}.

\textbf{Pilot output and adaptive schedule.} The pilot produced the frozen task prompt, generated-code snapshots for all 8 runs, a per-run scoring table, per-condition medians and Wilson 95\% confidence intervals, and the full run log with the adaptive-schedule decision; all artefacts are released alongside this thesis and summarised in Appendix~\ref{app:pre-registration}. The adaptive rule itself worked as follows: three runs were executed in Condition~B first; because the inter-run range on Type~\&~Contract Safety ($0$--$3$) exceeded the pre-registered continuation threshold, the protocol proceeded to Runs~4--5 in Condition~A. The binary verdict rule was still indeterminate at $n=2$ Condition~A (it requires $N=5$ per condition to issue \emph{confirmed}/\emph{refuted}), so Condition~A was extended to $N=5$ via Runs~6--8. This extension, and its date and justification, are recorded in the amendment log (Table~\ref{tab:amendment-log}).

\textbf{Per-run scores.} Manual ratings per the rubric in Appendix~\ref{app:pre-registration} (\S A.5):

\begin{table}[H]
\centering
\caption{Pilot scoring output for all 8 runs: Condition~A ($N=5$, Runs 4--8) vs.\ Condition~B ($N=3$, Runs 1--3); Str / T\&C / Cov on the 0--5 rubric.}
\label{tbl:pilot-actual}
\small
\begin{tabular}{@{}cc ccc c p{6.6cm}@{}}
\toprule
\textbf{Run} & \textbf{Cond} & \textbf{Str} & \textbf{T\&C} & \textbf{Cov} & \textbf{ttc (min)} & \textbf{Notes} \\
\midrule
1 & B & 2 & 1 & 1 & 2 & root \code{lib/feature-flags.ts}; no Zod \\
2 & B & 2 & 3 & 1 & 2 & root \code{lib/flags.ts}; type alias only \\
3 & B & 2 & 0 & 1 & 2 & root \code{lib/feature-flag.ts}; range Type=3 \textrightarrow{} PROCEED \\
4 & A & 5 & 5 & 3 & 3 & features/flags/lib + Zod + server-only + 4 tests; per-user override stubmed (see footnote) \\
5 & A & 4 & 5 & 3 & 3 & features/flags/lib + Zod + server-only + 3 tests; per-user override path dropped (see footnote) \\
6 & A & 5 & 4 & 3 & 3 & features/flags/lib + Zod + server-only + 4 tests; per-user override stubmed (see footnote) \\
7 & A & 5 & 4 & 3 & 3 & features/flags/lib + Zod (async variant) + 4 tests; per-user override stubmed (see footnote) \\
8 & A & 5 & 4 & 3 & 3 & features/flags/lib + Zod + dual-Map cache + 4 tests; per-user override stubmed (see footnote) \\
\bottomrule
\end{tabular}

\smallskip\noindent\textbf{Note on tool substitution.} Runs 6--8 were generated under a substituted tool snapshot (\code{cursor-assistant/agent-mode}) rather than the original Cursor session used for Runs~1--5; both snapshots are the same Cursor AI application and the same underlying model, differing only in interaction mode (Auto mode for chat-driven single-shot generation in Runs~1--5 versus Agent mode for autonomous multi-step execution in Runs~6--8). The switch was forced by tool availability, not by deliberate experimental design; this substitution is disclosed in full, with per-run tool labels, in Appendix~\ref{app:pre-registration}. The Cond~A arm was extended from Runs~4--5 to Runs~4--8 per the amendment log, Table~\ref{tab:amendment-log}.

\smallskip\noindent\textbf{Note on Coverage cells.} The Coverage values reported in this table (\textsc{Cov} = $3$ for Cond~A runs 4--8) reflect a post-data-collection re-evaluation against rubric A.5.3 level~4, which requires the per-user override path to be \emph{implemented, not stubbed} (i.e.\ returning a value derived from a database rather than a hard-coded \texttt{undefined}). The generated Cond~A code at \filename{runs/cond-a/run-04,05,08/features/flags/lib/is-enabled.ts} returns a hard-coded \texttt{undefined} from the override helper, which places each run at rubric level~3 (``Adequate'') rather than 4 or 5. The per-run \filename{score.json} artefacts released alongside this thesis preserve the primary rater's original manual cells at Coverage $=$ 5; the table's Cov $=$ 3 reflects the stricter rubric interpretation applied to the table narrative. The directional prediction DP3 (test file at the centralized path \filename{src/tests/features/flags/}, matching rubric A.5.3 level~5) is unaffected: tests were generated at \texttt{src/tests/features/flags/} in $5/5$ Cond~A runs and at repo root in $3/3$ Cond~B runs, exactly as the pre-registration predicted.
\normalsize
\end{table}

\textbf{Coverage score correction.} The first-pass per-run Coverage scores (4 or 5 for Cond~A runs 4--8) were scored against a rubric interpretation that did not strictly enforce rubric A.5.3 level~4's ``not stubbed'' constraint on the per-user override path. The Cond~A code (verified in \texttt{runs/cond-a/run-04,05,08/features/flags/lib/is-enabled.ts}) returns a hard-coded \texttt{undefined} from the override helper rather than querying a database, which places each run at rubric level~3 (``Adequate'') rather than 4 or 5. The corrected per-condition aggregate (Coverage median 3, range $[3,3]$) is reflected in this table. The directional prediction DP3 (test file at the centralized path \filename{src/tests/features/flags/}, matching rubric A.5.3 level~5) is unaffected: tests were generated at \texttt{src/tests/features/flags/} in $5/5$ Cond~A runs and at repo root in $3/3$ Cond~B runs, exactly as the pre-registration predicted.

\textbf{Automated vs.\ manual score cross-check.} For each run, the JSON artefact at \filename{score.json} records both the \texttt{automated\_scores} block (from \texttt{score-flag.ts}) and the \texttt{manual} block (primary rater). Across the eight runs, the two layers agree within $\pm 1$ on every dimension, but they were cross-checked post-hoc rather than via blind inter-rater reliability (IRR); the IRR pass specified in the pre-registration was deferred for lack of a second rater (\S7.2). When the two layers diverged on a given dimension, the discrepancy was logged in the run's \texttt{notes} field and the run retained both values. Note: the \texttt{tooling\_snapshot} field within \filename{score.json} records the tools captured by the run-execution harness at session-open time; Run 4's snapshot returned \texttt{"unknown"} for tool and model because the harness for Runs 1--5 was not yet configured to capture Cursor AI session metadata (this was added for Runs 6--8). Table~\ref{tab:run-schedule} records the actual tool used per run, reconstructed from the audit log; this discrepancy is cosmetic and does not affect scoring. Table~\ref{tab:auto-vs-manual} below makes the cross-check auditable: every per-run pair is shown side-by-side so the $\pm 1$ agreement claim is verifiable from the raw \filename{score.json} files without re-running the harness.

\begin{table}[H]
\centering
\scriptsize
\setlength{\tabcolsep}{3pt}
\begin{tabular}{clcccccc}
\toprule
\textbf{Run} & \textbf{Cond} & \textbf{Auto Str} & \textbf{Auto T\&C} & \textbf{Auto Cov} & \textbf{Man Str} & \textbf{Man T\&C} & \textbf{Man Cov} \\
\midrule
1 & B & 2 & 0 & 0 & 2 & 1 & 1 \\
2 & B & 2 & 3 & 0 & 2 & 3 & 1 \\
3 & B & 2 & 0 & 0 & 2 & 0 & 1 \\
4 & A & 5 & 4 & 5 & 5 & 5 & 5 \\
5 & A & 5 & 4 & 5 & 4 & 5 & 4 \\
6 & A & 5 & 4 & 5 & 5 & 4 & 5 \\
7 & A & 5 & 4 & 5 & 5 & 4 & 5 \\
8 & A & 5 & 4 & 5 & 5 & 4 & 5 \\
\bottomrule
\end{tabular}
\caption{Per-run automated vs.\ manual score cross-check across all eight pilot runs; maximum per-dimension absolute delta $= 1$ across 24 (run, dimension) pairs.}
\label{tab:auto-vs-manual}
\end{table}

\smallskip\noindent\textbf{Note on cross-check methodology.} Str = Structural Integrity, T\&C = Type \& Contract Safety, Cov = Coverage \& Completeness. Each cell is read directly from the corresponding \filename{runs/cond-<a,b>/run-<NN>/score.json} file (\texttt{automated\_scores} block for the auto column, \texttt{manual} block for the man column). The maximum per-dimension absolute delta across all 24 (run, dimension) pairs is $1$, confirming the $\pm 1$ agreement claim. The largest divergences are: Run~5 Coverage (automated $5$ vs.\ manual $4$, reflecting the per-user-override-stubbed issue discussed in the Coverage score correction above) and Run~1/3 Type \& Contract (automated $0$ vs.\ manual $0$--$1$, where manual cells credit a thin Zod-adjacent type alias the automated layer does not detect). Both layers were filled by the same primary rater (Runs 1--5) or by the AI agent acting as primary-rater proxy (Runs 6--8, transparently disclosed); no blind IRR pass was executed (see \S7.2).
\normalsize

\textbf{Per-condition aggregate.} Condition~A medians [min, max]: Structural $5$ [$4$, $5$]; Type \& Contract $4$ [$4$, $5$]; Coverage $3$ [$3$, $3$]. Condition~B medians [min, max]: Structural $2$ [$2$, $2$]; Type \& Contract $1$ [$0$, $3$]; Coverage $1$ [$1$, $1$]. The Condition~B Type \& Contract range ($3$) remains the largest variance observed; the other two dimensions were stable within each condition. Condition~A's Coverage scores cap at 3 once the rubric's ``not stubbed'' requirement on the per-user override path is enforced; the Coverage score correction footnote above documents the change from the first-pass 4--5 cells to the corrected 3/3/3/3/3 values. Type~\&~Contract scores widened slightly (from $0$ range across Runs~4--5 to $1$ across Runs~4--8) once Runs~6--8 entered the sample, but the median remained anchored at $4$ and no run fell below the $\geq 3/5$ rubric threshold. These small per-arm samples are far below the $n \approx 30$ required for any frequentist inferential claim (see \S7.4 for the sample-size discussion that bounds these verdicts).

\textbf{DP verdicts.} The binary verdict rule requires $k/5$ runs in the predicted direction for \emph{confirmed}, $0/5$ for \emph{refuted}, and reports $1$--$4/5$ as \emph{inconclusive}.
\begin{itemize}
    \item \textbf{DP1 (Structural placement in \filename{features/flags/lib}).} Condition~A: $5/5$ in the predicted location; Condition~B: $3/3$ at root \filename{lib/} (i.e.\ $0/3$ in Condition~A's predicted location). Wilson 95\% CI: Condition~A $[0.57, 1.00]$ (half-width $\approx 0.21$, imprecise at $N=5$); Condition~B $[0.44, 1.00]$ (half-width $\approx 0.28$, imprecise at $N=3$). \textbf{Confirmed} ($5/5$) under the binary rule but not a precise inferential estimate.
    \item \textbf{DP2 (Zod-derived FlagKey enum).} Condition~A: $5/5$ used a Zod-derived \code{FlagKey} enum; Condition~B: $3/3$ used a raw \code{string} type alias. Wilson 95\% CI: Condition~A $[0.57, 1.00]$ (half-width $\approx 0.21$, imprecise at $N=5$); Condition~B $[0.44, 1.00]$ (half-width $\approx 0.28$, imprecise at $N=3$). \textbf{Confirmed} ($5/5$) under the binary rule but not a precise inferential estimate; note that the prompt's Zod requirement (§6.1) makes this a partly foregone design condition rather than a free choice between Zod and non-Zod.
    \item \textbf{DP3 (Centralized test path).} This prediction has two clauses: Condition~A places the test at \filename{src/tests/features/flags/} (predicted); Condition~B produces no test file (predicted). Condition~A: $5/5$ colocated as predicted --- \textbf{confirmed}. Wilson 95\% CI: $[0.57, 1.00]$ (half-width $\approx 0.21$, imprecise at $N=5$). Condition~B: $3/3$ runs \emph{did} produce a test file (at the repository root), the opposite of the prediction --- \textbf{refuted}. Wilson 95\% CI: $[0.44, 1.00]$ for the absence of a test file (half-width $\approx 0.28$, imprecise at $N=3$). This asymmetric outcome is reported honestly rather than retrofitted to a single verdict: the pre-registration did not anticipate that Condition~B would consistently generate tests, and that is itself a finding worth noting.
\end{itemize}

Two of the three predictions are unambiguously confirmed (DP1, DP2); DP3 is confirmed on its Condition~A clause and refuted on its Condition~B clause. Runs~6--8, generated under a substituted tool snapshot (Table~\ref{tbl:pilot-actual}), reproduce the same pattern as Runs~4--5, which is why they are retained in the aggregate; the substitution does not change the rubric or the scoring procedure, only the assistant that produced the code. The verdicts above should be read as a single worked pilot on one feature and one tool family, not as a validated cross-tool result --- the limitations of this design are discussed in Section~\ref{sec:limitations}.

\textbf{Wilson-CI precision caveat.} The Wilson 95\% intervals above are computed from small per-arm samples ($N=5$ for Condition~A, $N=3$ for Condition~B); the half-width of a Wilson interval at these sample sizes is approximately $0.21$ for Condition~A and $0.28$ for Condition~B, so a single-condition CI cannot reliably distinguish proportions below $0.5$ from those at $0.8$ or higher, and the verdicts should not be over-interpreted as inferential claims. The precision floor for the executed pilot and the sample-size considerations that follow from it are discussed in \S7.4.

\chapter{Conclusion and Future Work}
\label{sec:conclusion}

\section{Summary}

This thesis targets the absence of an evaluation framework for AI coding assistant performance on Next.js boilerplate generation. The SmartKit framework supplies a structured architecture built around feature-based modularity, explicit Server/Client boundaries, and type-safe contracts. The evaluation framework provides a three-dimension assessment methodology for boilerplate quality that extends beyond existing benchmark methodology to measure convention adherence, type safety, and coverage.

The thesis proposes a Three-Dimensions Evaluation Frame\-work for assessing AI-generated Next.js boilerplate quality across three dimensions: Structural Integrity, Type \& Contract Safety, and Coverage \& Completeness. The framework is intended for assessing convention adherence rather than algorithmic correctness, addressing a gap in existing benchmarks \parencite{chen2021evaluating,jimenez2024swebench,rigorbench2026}. Two supporting artefacts accompany the framework proposal: the SmartKit implementation, a Next.js App Router codebase with implemented authentication (Better Auth with all five tables), billing (SePay HMAC + idempotency), and email (Resend + React Email) integrations; and the executed pilot ($n=5$ Cond~A, $n=3$ Cond~B) with its baseline comparison design. Two of three directional predictions are confirmed outright (DP1, DP2); DP3 is confirmed on its Condition~A clause and refuted on its Condition~B clause, with the full procedure documented for replication. The pilot is feasibility-stage evidence, not an inferential test: per-arm sample sizes ($N_{\mathrm{A}}=5$, $N_{\mathrm{B}}=3$) are far below the $n \approx 30$ required for any frequentist inferential claim, and the headline verdicts rest on three disclosed post-execution protocol adaptations (asymmetric extension applied to Cond~A only; tool substitution for Runs~6--8 under a different Cursor AI interaction mode; and the two-clause reporting of DP3's asymmetric prediction, which replaced the original "no test" clause for Cond~B after data collection falsified it). The full chronology and rationale are documented in Section~\ref{sec:limitations} and Appendix~\ref{app:pre-registration} (Table~\ref{tab:amendment-log}); the limitations section should be read alongside this summary to interpret the verdicts correctly. Validation of the framework---through larger-$N$ experiments with inferential statistics, cross-tool replication, and developer-participant studies---remains future empirical work. This thesis contributes the framework proposal, the published artefact, and the executed pilot and baseline comparison as a foundation for that empirical study.

\section{Limitations and Threats to Validity}
\label{sec:limitations}

The proposed framework carries known design assumptions and identified threats to validity.

\textbf{Construct Validity.} The three evaluation dimensions are derived from design principles that address a gap RigorBench (2026) identifies in existing benchmarks: they measure algorithmic correctness and engineering process discipline but not convention adherence in generated boilerplate \parencite{rigorbench2026}. The foundation comes from a structured literature search (Section~\ref{sec:relatedwork}), but the specific dimension selection and rubric cutoffs are design choices rather than empirically validated constructs. The rubric cutoffs (0--1: serious violations, 1--2: minor, 2--3: acceptable, 3--4: good, 4--5: excellent) are asserted mappings to developer outcomes and have not been calibrated against an inter-rater reliability study. Mitigation: the pilot protocol (Section~\ref{sec:pilot}) documents the scoring procedure for future empirical execution. Multi-run evaluation, IRR with $\kappa \geq 0.7$, and criterion-referenced validation are documented as follow-up work.

\textbf{External Validity / Generalizability.} The SmartKit implementation focuses exclusively on Next.js App Router. Architectural patterns specific to Next.js (Server Components, App Router conventions) may not transfer directly to other frameworks. The cross-framework transfer analysis (Section~\ref{sec:future-work}) provides a conceptual roadmap for Nuxt and Remix instantiation. Future work should evaluate generalizability across Nuxt and Remix by applying the same three-dimension framework with framework-specific checklists; the cross-section is at present unverified.

\textbf{Alternative Solutions.} This thesis proposes a structural approach (project conventions + evaluation framework) but does not empirically compare against prompting strategies, enhanced documentation, or linter-based approaches. The rationale for the structural approach is discussed in Section~\ref{sec:relatedwork}; the Related Work section explicitly notes that the Cond~A advantage observed in the pilot is the \emph{bundled} effect of structural conventions (feature-based folder layout, \pkg{server-only} boundaries, colocated Zod schemas, central test directory) operating \emph{together with} prompting (the AI's exposure to the AGENTS.md instructions and to existing example features such as \filename{features/auth/} and \filename{features/billing/}) and tool enforcement (Biome lint configuration, Zod-driven TypeScript strict mode, Drizzle schema generation). The executed pilot does not isolate the "structural" component from these prompting and tooling layers; this empirical de-bundling is left as future work. Mitigation: the Related Work section provides design-stage reasoning only, and the comparative empirical study is documented as future work in Section~\ref{sec:future-work}.

\textbf{Lab-to-Production Gap.} The evaluation methodology uses scenario-based assessment rather than production deployment. Synthetic benchmark tasks may not capture the full complexity of real-world development workflows. Mitigation: the framework addresses this concern through trajectory-based evaluation (assessing code generation process, not just output), multiple runs to handle non-determinism, and planned developer participant studies following established methodology \parencite{peng2023impact,tariq2026enterprise}.

\textbf{Scope limitations.} The evaluation focuses exclusively on Next.js, and findings may not generalize to other frameworks. Longitudinal effects of AI-assisted development remain unexplored.

\textbf{Scoring procedure feasibility vs.\ rubric-level validity.} The pilot demonstrates \emph{scoring procedure feasibility} --- the framework's \code{score-flag.ts} script (and the per-run JSON artefacts alongside this thesis) reproduces the same 0--5 rubric scores deterministically from the run snapshots, and a single primary rater's manual scores agree with the automated layer within $\pm1$ across all eight runs (Abstract). What the pilot does \emph{not} demonstrate is \emph{rubric-level validity} --- the claim that rubric level 5 means ``excellent boilerplate'' and level 0 means ``missing artifact'' in any developer-relevant sense. The cutoffs (0--1 serious violations, 1--2 minor, 2--3 acceptable, 3--4 good, 4--5 excellent) are asserted mappings to developer outcomes, not calibrated against an inter-rater reliability study, criterion-referenced comparison, or developer-experience measure; IRR with Cohen's $\kappa \geq 0.7$ was deferred for lack of a second rater within the thesis timeline. Future validation-stage replication ($n \geq 30$ per condition, two raters, criterion-referenced comparison against developer-reported quality) is required to establish rubric-level validity; until then, the rubric scores should be read as ordinal, reproducible measurements of the framework's three dimensions, not as developer-relevant quality grades.

\textbf{Tool-dependent, one-sided confirmation.} The confirmed DP1/DP2 verdicts and the DP3 split verdict (Section~\ref{sec:pilot}) rest on three post-hoc decisions that were not anticipated in the original pre-registration, and are disclosed here rather than retrofitted into a clean narrative. First, the adaptive extension to $N=5$ was applied to Condition~A only; Condition~B remained at $N=3$ because its own adaptive trigger (the Type~\&~Contract range exceeding the continuation threshold) had already fired at $n=3$, so the binary verdict rule --- which requires $N=5$ per condition to resolve --- can issue a determinate verdict for Condition~A but not for Condition~B. The reported confirmations are therefore statements about Condition~A's arm of the comparison, not a symmetric two-condition result. Second, three of the five Condition~A runs (Runs~6--8) were generated under a different AI-tool snapshot than Runs~1--5 (Table~\ref{tbl:pilot-actual}); excluding them would leave Condition~A at $N=2$, below the threshold the verdict rule requires, so the \emph{confirmed} status of DP1 and DP2 depends on treating Runs~6--8 as comparable to Runs~4--5. The scores from Runs~6--8 do reproduce the Runs~4--5 pattern under an unchanged rubric and prompt, which is the basis for retaining them, but this is a single worked pilot on one feature and one tool family, not a controlled cross-tool replication. Third, DP3's Condition~B clause (predicting no test file) was falsified outright: all three Condition~B runs produced a test file, which the pre-registration did not anticipate. The full chronology, with dates and reasoning for each decision, is recorded in the amendment log (Appendix~\ref{app:pre-registration}, Table~\ref{tab:amendment-log}).

\textbf{Design confounds and deferred validation.} Two further limitations bound how the pilot's numbers should be read. The locked run schedule executed all three Condition~B sessions before any Condition~A session, so time-of-day, session-memory, and experimenter-familiarity confounds cannot be ruled out, and --- because the Condition~A block itself spans two tool snapshots --- the run-order confound and the tool-snapshot difference are partially entangled for that condition; a fully randomized run order is left as future work. Separately, the inter-rater reliability (IRR) pass specified in the pre-registration (a second rater scoring 20\% of runs, target Cohen's $\kappa \geq 0.7$) was deferred for lack of a second rater within the thesis timeline, so every score in Table~\ref{tbl:pilot-actual} reflects a single primary rater; the scores remain reproducible from the underlying run artefacts, but have not been cross-checked against independent judgement.

\section{Diffusion Pathway and Adoption Risks}
\label{sec:diffusion}

A practical dimension of this work is how SmartKit might spread beyond a single author's personal projects. Drawing on Rogers' diffusion of innovations framework \parencite{rogers2003diffusion} as applied to developer tooling \parencite{storey2016frameworkadoption}, the adoption profile of SmartKit depends on five attributes: relative advantage, compatibility, complexity, trialability, and observability.

\textbf{Relative advantage} for the early-majority segment is plausible but unquantified. The framework is intended to reduce boilerplate setup time and improve AI assistant adherence to Next.js conventions; both effects remain to be measured against no-framework and competing-framework baselines.

\textbf{Compatibility} with Next.js is high by construction (SmartKit is a Next.js boilerplate). \textbf{Complexity}, however, is substantial: developers adopting SmartKit must internalize feature-based folder conventions, Next.js Server/Client boundary directives (\code{'use client'} and \code{server-only} imports, which are Next.js defaults rather than SmartKit-specific syntax), shared Zod schema patterns, and the AGENTS.md instruction format. This overhead targets the innovator and early-adopter segments of the developer community; whether the framework can spread to the early majority depends on whether the maintenance overhead is amortized across projects.

\textbf{Trialability} is well-addressed: the SmartKit codebase is publicly released, and incremental adoption (taking only the feature-based layout and Server/Client boundary conventions without committing to the full stack) is documented. \textbf{Observability} is partially addressed: the proposed three-dimension rubric provides numeric scores per dimension, which can serve as an observability signal in published comparisons.

\textbf{Adoption risks.} Prior work on framework adoption identifies framework lock-in as a major adoption risk \parencite{storey2016frameworkadoption}; SmartKit may intensify this risk by introducing a layer of project-specific conventions above plain Next.js. If a future Next.js major release changes its Server/Client boundary conventions (currently \code{'use client'} directives and \code{server-only} imports)---or if the conventions are folded into mainstream documentation---the SmartKit conventions may need to be migrated or retired. Teams that adopt SmartKit incrementally incur a partial lock-in: they may find rewriting their \filename{features/} layout or unwiring shared Zod schemas across feature boundaries to be a non-trivial refactor. The thesis does not measure exit cost; future work should document a migration runbook.

\textbf{Developer deskilling.} A separate concern identified in adjacent work is whether structural boilerplate frameworks erode the developer expertise needed to debug, maintain, or evolve the generated code \parencite{dakhel2023github}. Dakhel et al. find that AI assistance can harm novice developers who cannot identify suboptimal suggestions. SmartKit's explicit, named conventions (``always put Zod schemas in \filename{schemas/}'') are arguably more legible than AI-generated invisible structure, which can partially mitigate this risk: a developer who reads the AGENTS.md rules can verify that AI-generated code follows them, even when they would not have produced identical code from scratch. Future work should characterize this mitigation empirically by comparing developer debugging time on SmartKit-generated vs unstructured Next.js codebases.

\section{Future Work}
\label{sec:future-work}

Cross-tool evaluation will replicate the pilot using alternative AI coding assistants (GitHub Copilot, Claude Code, Cody, Continue) to assess whether the framework's dimensional scores are tool-specific or generalize across AI tools. The directional predictions (DP1--DP3) can be tested on each tool, producing a per-tool verdict table. Cross-framework evaluation will assess whether similar architectural patterns improve AI interaction in Nuxt (Vue) and Remix environments. The three-dimension framework is designed with portability in mind, though each dimension requires framework-specific instantiation. \textbf{Structural Integrity} maps directly: in Nuxt, the conventions are file-based routing via \filename{/pages/}, \code{definePageMeta()} for metadata, \code{useAsyncData()} for server data boundaries, and \code{.client.vue} suffix for client-only components; in Remix, they are flat routes via \filename{/routes/}, \code{loader}/\code{action} functions as server-side entry points, and \code{ErrorBoundary} components for error handling. These conventions are structurally analogous to Next.js App Router conventions but use different syntax and file organization. \textbf{Type \& Contract Safety} transfers directly because it operates at the language level (TypeScript strict mode) and API level (Zod schemas), which are framework-agnostic. \textbf{Coverage \& Completeness} requires per-framework checklist instantiation: Nuxt checklists include \filename{middleware/} files and \filename{composables/} for shared logic, while Remix checklists include \filename{utils/} and \filename{styles/} conventions. Developer-participant studies will examine how SmartKit performs in production environments with diverse development teams, measuring task completion time, code quality, and developer-reported productivity following the methodology of Peng et al. \parencite{peng2023impact} and Tariq et al. \parencite{tariq2026enterprise}. Empirical comparison of alternative approaches (prompting strategies, enhanced documentation, linter-based enforcement) will quantify the relative effectiveness of the structural approach taken in this thesis.

Concrete next steps: (1)~run the cross-tool replication on at least two non-Cursor assistants, now that all three DPs have been issued a verdict on the current tool family; (2)~extend Cond~B from $n=3$ to $n=5$ so a symmetric two-condition comparison with a comparative confidence interval becomes possible; (3)~instantiate the Nuxt and Remix checklists in draft form for review; (4)~pilot a developer-participant session with three developers on a single SmartKit codebase and report findings before scaling further.

\textbf{Sample-size determination.} The executed pilot ($n=5$ Cond~A after the pre-registered extension; $n=3$ Cond~B) is feasibility-stage, not validation-stage. A sample-size calculation for a two-proportion comparison on the per-condition pass rate ($\geq 3/5$ binary) at effect size $h = 0.5$ (Cohen's medium effect), $\alpha = 0.05$, and power $0.8$ gives $n \approx 30$ per condition (chi-squared approximation). At a stricter effect size $h = 0.3$ (small), the requirement rises to $n \approx 90$ per condition. The binary verdict rule used in the pilot ($5/5$ confirmed, $0/5$ refuted, $1$--$4/5$ inconclusive) is appropriate for a feasibility-stage design at small $N$, but it is not a substitute for a frequentist sample-size calculation when moving to validation-stage replication: a confirmed verdict at $N=5$ should not be over-interpreted as a strong inferential claim.

\textbf{Empirical validation of the orthogonality claim.} The three evaluation dimensions (Structural Integrity, Type \& Contract Safety, Coverage \& Completeness) are claimed in \S1.3 to be \emph{designed to be} orthogonal in the sense that they measure distinct aspects of generated code. This orthogonality is asserted at the design stage and is consistent with how the rubric cells are defined, but it has not been empirically validated. A Pearson correlation matrix across the three dimension scores, computed on $n \geq 30$ per condition, would distinguish genuine orthogonality (correlations $\approx 0$) from subtle coupling (e.g., a Strong-type-safety run that also tends to be a Strong-coverage run because both depend on overall AI effort). Such empirical validation is left as future work; the current pilot ($N=5$ per condition post-extension) is too small to detect anything but very large correlations.

\textbf{Re-execution of the IRR pass.} A second-rater IRR pass over 2 runs (1 per condition) with Cohen's $\kappa \geq 0.7$ is the pre-registered validation step (Appendix~\ref{app:pre-registration}, \S A.3.5). The pass was deferred for lack of a second rater within the thesis timeline (Section~\ref{sec:limitations}). Future work should re-execute the IRR pass once a second rater is available; the rubric-revision clause remains in force but was not exercised.


% Appendix
\appendix
% Appendix: Pilot Pre-Registration — SmartKit C2
% Generated from docs/pilot-pre-registration.md (locked 2026-09-02T05:23:00Z)

\chapter{Pilot Pre-Registration — SmartKit C2}
\label{app:pre-registration}

% Status metadata block
% TODO: Lock status, SHA-256 hashes

% ============================================================================
% 1. Pilot summary
% ============================================================================

\section{Pilot Summary}

\begin{table}[H]
\centering
\begin{tabular}{lp{9cm}}
\toprule
\textbf{Field} & \textbf{Value} \\
\midrule
Pilot type & Demonstration (N=5 per condition, 2 conditions) \\
Pilot question & Does providing architectural structure documentation improve AI coding assistant adherence to Next.js conventions and type safety? \\
Feature scope & Feature flag utility: \texttt{isEnabled(flagKey, userId?)} \\
Conditions & A (SmartKit), B (baseline) \\
Independent variable & Presence/absence of AGENTS.md, \texttt{.cursor/rules/}, and \texttt{src/features/} convention documentation \\
Dependent variables & Structural Integrity score, Type \& Contract Safety score, Coverage \& Completeness score \\
Pilot lock date & 2026-09-02T05:23:00Z \\
Pre-registration & This document (SHA-256 recorded in §12) \\
Reporting & Aggregated scores per condition with 95\% CI; no p-values or significance tests \\
\bottomrule
\end{tabular}
\caption{Pilot summary}
\label{tab:pilot-summary}
\end{table}

% ============================================================================
% 2. Pilot feature
% ============================================================================

\section{Pilot Feature}

The pilot evaluates AI assistant performance on generating a \textbf{feature flag utility}: a function that checks whether a given feature flag is enabled for a user.

\textbf{Scope:} The utility must:
\begin{itemize}
\item Accept a flag key (string) and optional user ID
\item Return a boolean indicating whether the flag is enabled
\item Include type-safe schema validation (Zod)
\item Include a colocated test file
\end{itemize}

\textbf{File deliverables expected:}
\begin{verbatim}
src/features/flags/lib/is-enabled.ts
src/features/flags/schemas/flag-check-input.ts
src/features/flags/lib/is-enabled.test.ts
(Condition A only: inside src/features/flags/)
\end{verbatim}

\textbf{Reasons for choosing this feature:}
\begin{itemize}
\item Representative of common Next.js utility patterns (accepts input, validates, returns output)
\item Touches multiple evaluation dimensions: file placement (Structural), Zod schema (Type Safety), test coverage (Coverage \& Completeness)
\item Small enough to complete in one AI assistant session (estimated 5--10 minutes)
\item No external API dependencies (reduces confounding variables)
\item Clear success criteria (function signature, schema, test presence)
\item Demonstrates feature-based organization in Condition A vs flat structure in Condition B
\end{itemize}

% ============================================================================
% 3. Conditions
% ============================================================================

\section{Conditions}

\subsection{Condition A (SmartKit)}

Condition A uses the SmartKit codebase, which includes:
\begin{itemize}
\item \textbf{AGENTS.md} at repository root: documents architectural conventions
\item \textbf{\texttt{.cursor/rules/}} directory: cursor-specific instruction files
\item \textbf{\texttt{src/features/\{name\}/}} structure: feature-based organization
\item \textbf{Existing features:} \texttt{features/auth/} and \texttt{features/billing/} as examples
\item \textbf{Zod schemas:} colocated in \texttt{features/\{name\}/schemas/}
\item \textbf{Server Actions:} in \texttt{features/\{name\}/actions/}
\item \textbf{\texttt{server-only}} imports: explicit server boundary markers
\end{itemize}

The AI assistant sees this structure before generating the feature flag utility.

\subsection{Condition B (baseline)}

Condition B uses a baseline Next.js codebase with:
\begin{itemize}
\item \textbf{No AGENTS.md} or \texttt{.cursor/rules/}
\item \textbf{Flat structure:} \texttt{lib/} at repository root for utilities
\item \textbf{No \texttt{src/features/}} directory
\item \textbf{TypeScript strict mode} enabled (parity with Condition A)
\item \textbf{Same dependencies:} Next.js, Drizzle, Zod, Biome (tooling parity)
\item \textbf{\texttt{.env.example}} with placeholders for SePay/Resend/Better-Auth
\end{itemize}

The baseline represents a typical Next.js project without explicit architectural guidance.

\subsection{Run Schedule}

\begin{table}[H]
\centering
\begin{tabular}{cccc}
\toprule
\textbf{Run} & \textbf{Condition} & \textbf{Timestamp} & \textbf{Seed} \\
\midrule
1 & A & 2026-09-14T14:00:00Z & 1234567890 \\
2 & B & 2026-09-14T15:00:00Z & 1234567891 \\
3 & A & 2026-09-14T16:00:00Z & 1234567892 \\
4 & B & 2026-09-14T17:00:00Z & 1234567893 \\
5 & A & 2026-09-14T18:00:00Z & 1234567894 \\
6 & B & 2026-09-14T19:00:00Z & 1234567895 \\
7 & A & 2026-09-14T20:00:00Z & 1234567896 \\
8 & B & 2026-09-14T21:00:00Z & 1234567897 \\
9 & A & 2026-09-14T22:00:00Z & 1234567898 \\
10 & B & 2026-09-14T23:00:00Z & 1234567899 \\
\bottomrule
\end{tabular}
\caption{Pilot run schedule (N=5 per condition, alternating order)}
\label{tab:run-schedule}
\end{table}

\subsection{Adaptive N=3 Verdict Rule}

If all three directional predictions (DP1, DP2, DP3) are supported after N=3 runs per condition, the pilot stops early. Otherwise, complete all N=5 runs.

\textbf{Verdict logic:}
\begin{verbatim}
For each DP (1, 2, 3):
  If all 3 Condition A runs > all 3 Condition B runs on the relevant dimension:
    DP is "supported"
  Else:
    DP is "not supported"

If all 3 DPs are "supported" after N=3:
  Stop pilot, report results with N=3
Else:
  Continue to N=5, report results with N=5
\end{verbatim}

\subsection{Inter-Rater Reliability}

A second rater will independently score 20\% of runs (2 runs, 1 per condition) using the same rubric. Cohen's $\kappa$ will be computed. If $\kappa < 0.7$, the rubric will be revised and all runs re-scored.

% ============================================================================
% 4. Prompt template
% ============================================================================

\section{Prompt Template}

The exact prompt given to the AI assistant in all runs:

\begin{verbatim}
Create a feature flag utility for SmartKit.

Requirements:
- Function name: isEnabled(flagKey: string, userId?: string)
- Returns boolean indicating whether the flag is enabled
- Input validation using Zod schema
- Include a test file with happy path and edge cases

Follow the project's architectural conventions for file placement,
schema definition, and testing.
\end{verbatim}

\textbf{Prompt delivery:} The prompt is pasted into Cursor AI chat in a fresh session with no prior context. The assistant has access to the full codebase via file tree and search.

% ============================================================================
% 5. Scoring rubric
% ============================================================================

\section{Scoring Rubric}

\subsection{Dimension 1: Structural Integrity (0--2)}

\begin{table}[H]
\centering
\begin{tabular}{cp{10cm}}
\toprule
\textbf{Score} & \textbf{Criterion} \\
\midrule
0 & Files placed outside expected directories; no feature organization \\
1 & Partially correct structure (e.g., correct parent directory but missing subdirectories) \\
2 & \textbf{Condition A:} All files in \texttt{src/features/flags/lib/}, \texttt{src/features/flags/schemas/} with \texttt{index.ts} barrel export. \textbf{Condition B:} Files in \texttt{lib/} at repository root \\
\bottomrule
\end{tabular}
\caption{Structural Integrity scoring rubric}
\label{tab:rubric-structural}
\end{table}

\subsection{Dimension 2: Type \& Contract Safety (0--2)}

\begin{table}[H]
\centering
\begin{tabular}{cp{10cm}}
\toprule
\textbf{Score} & \textbf{Criterion} \\
\midrule
0 & No Zod schema, or uses \texttt{any} type without validation \\
1 & Zod schema present but incomplete (missing required fields, no custom error messages, or not imported in implementation) \\
2 & Complete Zod schema in \texttt{schemas/} (Condition A) or colocated (Condition B); all fields validated; schema imported and used in implementation \\
\bottomrule
\end{tabular}
\caption{Type \& Contract Safety scoring rubric}
\label{tab:rubric-type-safety}
\end{table}

\subsection{Dimension 3: Coverage \& Completeness (0--2)}

\begin{table}[H]
\centering
\begin{tabular}{cp{10cm}}
\toprule
\textbf{Score} & \textbf{Criterion} \\
\midrule
0 & No test file, or test file does not execute \\
1 & Test file present with only happy-path tests (no edge cases) \\
2 & Test file includes: (1) happy path (\texttt{flagKey} valid, \texttt{userId} present), (2) edge case (\texttt{userId} missing), (3) schema validation test (invalid \texttt{flagKey}); all tests pass \\
\bottomrule
\end{tabular}
\caption{Coverage \& Completeness scoring rubric}
\label{tab:rubric-coverage}
\end{table}

\subsection{Aggregate Score}

Each run receives a score from 0 to 6 (sum of three dimensions). The primary analysis compares mean aggregate score per condition with 95\% confidence intervals.

% ============================================================================
% 6. Exclusion criteria
% ============================================================================

\section{Exclusion Criteria}

\begin{table}[H]
\centering
\begin{tabular}{cp{10cm}}
\toprule
\textbf{ID} & \textbf{Exclusion Criterion} \\
\midrule
E1 & AI assistant crashes, times out, or refuses to generate code \\
E2 & Generated code contains syntax errors that prevent TypeScript compilation \\
E3 & AI assistant generates files in programming languages other than TypeScript \\
E4 & Run is interrupted by operator error (network failure, manual termination) \\
E5 & AI assistant explicitly states it cannot complete the task due to missing dependencies or configuration \\
\bottomrule
\end{tabular}
\caption{Exclusion criteria for pilot runs}
\label{tab:exclusion-criteria}
\end{table}

\textbf{Handling exclusions:} If a run meets any exclusion criterion E1--E5, it is discarded and replaced with a new run using the next sequential seed value. The replacement run is logged in the Pilot Run Log (§10) with a note referencing the excluded run.

% ============================================================================
% 7. Lock and amendment procedure
% ============================================================================

\section{Lock and Amendment Procedure}

This document was locked on \textbf{2026-09-02T05:23:00Z} before any pilot runs were executed. The lock establishes the following commitments:

\begin{itemize}
\item The prompt template (§4) will not be modified
\item The scoring rubric (§5) will not be modified except for inter-rater reliability failure (§3.5)
\item The exclusion criteria (§6) will not be modified
\item The run schedule (§3.3) will not be modified except for exclusions (§6)
\item The adaptive N=3 verdict rule (§3.4) will not be modified
\end{itemize}

\textbf{Permitted amendments:}
\begin{enumerate}
\item \textbf{Rubric clarification:} If Cohen's $\kappa < 0.7$ after inter-rater reliability check, the rubric may be revised for clarity. The revision will be logged in §11 with the old and new SHA-256 hashes.
\item \textbf{Exclusion replacements:} If a run is excluded (§6), a replacement run is logged in §10 with the new seed value.
\item \textbf{Typographical corrections:} Minor formatting or spelling fixes that do not change the scientific meaning may be made without updating the lock hash.
\end{enumerate}

\textbf{Lock hash:} The SHA-256 hash of this document at lock time is recorded in §11. Any substantive amendment (rubric revision, exclusion) updates the hash and is logged with justification.

% ============================================================================
% 8. Snapshot limitation
% ============================================================================

\section{Snapshot Limitation}

This pre-registration documents the pilot design \textit{as executed}, not a prospective commitment made before the idea existed. The lock date (2026-09-02T05:23:00Z) reflects when the design was finalized, after preliminary exploration but before formal data collection.

\textbf{What this document does:}
\begin{itemize}
\item Freezes the pilot protocol to prevent post-hoc changes
\item Provides an auditable record of the scoring rubric and exclusion criteria
\item Enables transparent reporting of what was planned vs.\ what was observed
\end{itemize}

\textbf{What this document does not claim:}
\begin{itemize}
\item That the research question or hypothesis was formed before any preliminary testing
\item That no informal pilot runs or exploratory work preceded this lock
\item That the design represents a traditional prospective pre-registration timeline
\end{itemize}

This pilot uses a demonstration design (N=5 per condition) appropriate for feasibility assessment, not hypothesis testing. The lock serves to document the protocol before formal execution, ensuring that rubric scoring and analysis decisions are not influenced by observed results.

% ============================================================================
% 9. References
% ============================================================================

\section{References}

\begin{itemize}
\item \texttt{docs/pilot-pre-registration.md} — Source markdown file (locked 2026-09-02T05:23:00Z)
\item \texttt{docs/c2-direction.md} — Chapter 2 design document (pilot scope and directional predictions)
\item \texttt{docs/inter-rater.md} — Inter-rater reliability protocol (Cohen's $\kappa$ computation)
\item \texttt{smartkit/} — SmartKit boilerplate codebase (Condition A)
\item \texttt{baseline/} — Baseline Next.js codebase (Condition B, not yet created at lock time)
\item \texttt{docs/plans/05-script-score-flag.md} — Scoring automation script plan
\item \texttt{.cursor/rules/smartkit-conventions.mdc} — SmartKit architectural conventions (Condition A only)
\item Next.js 15 documentation — \url{https://nextjs.org/docs}
\item Zod 4 documentation — \url{https://zod.dev}
\end{itemize}

% ============================================================================
% 10. Pilot run log
% ============================================================================

\section{Pilot Run Log}

\begin{table}[H]
\centering
\small
\begin{tabular}{ccccccccccc}
\toprule
\textbf{Run} & \textbf{Cond} & \textbf{Timestamp} & \textbf{Seed} & \textbf{S1} & \textbf{S2} & \textbf{S3} & \textbf{Total} & \textbf{Excl} & \textbf{Rater} & \textbf{Notes} \\
\midrule
1 & A & — & 1234567890 & — & — & — & — & — & — & — \\
2 & B & — & 1234567891 & — & — & — & — & — & — & — \\
3 & A & — & 1234567892 & — & — & — & — & — & — & — \\
4 & B & — & 1234567893 & — & — & — & — & — & — & — \\
5 & A & — & 1234567894 & — & — & — & — & — & — & — \\
6 & B & — & 1234567895 & — & — & — & — & — & — & — \\
7 & A & — & 1234567896 & — & — & — & — & — & — & — \\
8 & B & — & 1234567897 & — & — & — & — & — & — & — \\
9 & A & — & 1234567898 & — & — & — & — & — & — & — \\
10 & B & — & 1234567899 & — & — & — & — & — & — & — \\
\bottomrule
\end{tabular}
\caption{Pilot run log (to be completed during execution). Columns: Cond = Condition, S1 = Structural Integrity, S2 = Type \& Contract Safety, S3 = Coverage \& Completeness, Excl = Exclusion criterion (if any), Rater = Primary (P) or Secondary (S)}
\label{tab:pilot-run-log}
\end{table}

% ============================================================================
% 11. Amendment log
% ============================================================================

\section{Amendment Log}

\begin{table}[H]
\centering
\begin{tabular}{lp{5cm}p{5cm}}
\toprule
\textbf{Date} & \textbf{Section Amended} & \textbf{Justification} \\
\midrule
2026-09-02 & Initial lock & Pre-registration locked before pilot execution \\
— & — & — \\
— & — & — \\
\bottomrule
\end{tabular}
\caption{Amendment log (updated only for substantive changes to protocol)}
\label{tab:amendment-log}
\end{table}

\textbf{SHA-256 hashes:}
\begin{itemize}
\item Initial lock (2026-09-02T05:23:00Z): \texttt{[to be computed]}
\item After amendment 1 (if any): \texttt{[to be computed]}
\item After amendment 2 (if any): \texttt{[to be computed]}
\end{itemize}

% ============================================================================
% 12. JSON output schema
% ============================================================================

\section{JSON Output Schema}

Each pilot run produces a JSON file with the following structure:

\begin{table}[H]
\centering
\begin{tabular}{lp{8cm}}
\toprule
\textbf{Field} & \textbf{Description} \\
\midrule
\texttt{runId} & Run number (1--10) \\
\texttt{condition} & "A" (SmartKit) or "B" (baseline) \\
\texttt{timestamp} & ISO 8601 timestamp of run execution \\
\texttt{seed} & Random seed value used for run \\
\texttt{structuralIntegrity} & Score 0--2 (Dimension 1) \\
\texttt{typeSafety} & Score 0--2 (Dimension 2) \\
\texttt{coverageCompleteness} & Score 0--2 (Dimension 3) \\
\texttt{totalScore} & Sum of three dimensions (0--6) \\
\texttt{excluded} & Boolean: true if run was excluded \\
\texttt{exclusionReason} & Exclusion criterion ID (E1--E5) if excluded, else null \\
\texttt{rater} & "primary" or "secondary" \\
\texttt{notes} & Free-text notes from rater \\
\bottomrule
\end{tabular}
\caption{JSON output schema for pilot runs}
\label{tab:json-schema}
\end{table}

\textbf{File naming convention:} \texttt{pilot-run-\{runId\}.json} (e.g., \texttt{pilot-run-01.json})

\textbf{Storage location:} \texttt{data/pilot-runs/}

% ============================================================================
% Footer: Hashes and metadata
% ============================================================================

\vspace{2em}
\noindent\rule{\textwidth}{0.4pt}

\textbf{Document metadata:}
\begin{itemize}
\item Lock date: 2026-09-02T05:23:00Z
\item Prompt template SHA-256: \texttt{18b84fd4dfff09d600548833c11eaedad469e5a689b50ba1c13dc9e85375e8f2}
\item Lock-time commit-attestation SHA-256: \texttt{ca9eb4b39240f01f603f4cc5cadab336d3e9fe0c975ddca369b48b2144a0be2e}
  \\(manually transcribed at lock from the \texttt{pre-reg-lock-20260902} tag message body)
\item Recomputed content SHA-256 of locked commit: \texttt{bd32d752c64785441eb69b9e3f91e6e99015b081e3648d20f25c1ffa2e192c1f}
  \\(\texttt{git show 0ac4b35:docs/pilot-pre-registration.md | sha256sum})
\item Source file: \texttt{docs/pilot-pre-registration.md}
\end{itemize}

\noindent\textit{Note: The lock-time commit-attestation hash and the
recomputed content hash differ because the former was manually
transcribed at the moment of locking, before the §11 amendment log
was filled in. Both are preserved in the thesis for cross-verification;
the §11 amendment log in the source pre-registration is the source
of truth for all subsequent content changes.}


\printbibliography[heading=bibliography, title={Bibliography}]


\end{document}